\documentclass[11pt]{article}
\usepackage{amsmath,amssymb,amsthm,geometry}
\usepackage[hidelinks]{hyperref}
\geometry{margin=1in}
\newtheorem{proposition}{Proposition}
\newtheorem{remark}{Remark}
\newcommand{\R}{\mathbb R}
\newcommand{\dd}{\,\mathrm d}
\newcommand{\proj}{\operatorname{proj}}
\newcommand{\av}[1]{\left\langle #1\right\rangle}
\newcommand{\Cov}{\mathcal C}
\newcommand{\Cross}{\mathcal B}
\usepackage{mathrsfs}
\newtheorem{theorem}{Theorem}

\begin{document}
\section{The actual wave increment and the source's three compatibility targets}
\label{momop:section}

The source is the retained 166-page manuscript. Its SHA-256 is given
in these two consecutive segments:
\[
\begin{gathered}
\texttt{0e779481c4da40bd28d1e642e1d8ca57}\\
\texttt{447d129610df28dfa5a11e9af8ae228f}.
\end{gathered}
\]
The source operations used here are (8.3)--(8.27) and the complete
four-step cycle of Proposition 9.6, on printed pages 89--111.
Its coefficient classes are Definitions 6.4--6.5 and Proposition 6.6,
printed pages 69--72. The source estimates quoted below are attributed
to those statements. The following calculation of the actual candidate
targets and all pressure-induced changes is a direct finite calculation.
It does not infer source iteration exponents from finite smoothness.
The source correction cycle is its viscosity-one cycle. The physical
increment identities below separately retain any fixed
$\nu_{\rm NS}>0$; using another physical viscosity in those identities
does not identify the resulting candidate with the source's witness.

\subsection{The actual covariance increment before pressure reconstruction}

Use right-handed cylindrical coordinates $(r,\theta,z)$. An overbar
means the normalized angular average followed by normalized auxiliary
Haar average, taken at fixed slow variables before phase evaluation.
The old source state has actual zero-angular-mean wave $w_*$, including
all its curl corrections. The new candidate is the proved finite field
\[
 w_y=\sum_{j,a}y_a(t_{c,j})Z_{ja},\qquad
 \pi_y=\sum_{j,a}y_a(t_{c,j})\Pi_{ja},\qquad
 \langle w_y\rangle_\theta=\langle\pi_y\rangle_\theta=0.
\]
The $Z_{ja},\Pi_{ja}$ are the complete physical shapes from the
third-growth quantitative calculation; $y_a$ is the actual coupled
pair, with $t_{c,j}=t_b+(r_a/L_{s,j})v_j$. No amplitude factor is
removed from an auxiliary average: it can depend on $v_j(Y)$.
Write the exact double-averaged covariance increment as
\begin{equation}\label{momop:C}
 C_{ab}=
 \overline{w_{*,a}w_{y,b}+w_{y,a}w_{*,b}+w_{y,a}w_{y,b}},
 \qquad a,b\in\{r,\theta,z\}.
\end{equation}
Equivalently, substitute the displayed finite sum for both occurrences
of $w_y$ inside this average. In particular $C$ retains every old-wave
cross term and every same-label quadratic term. It is symmetric.
All its radial integrals below are finite because the new fields are
smooth and supported in a compact interior annulus. The old field
need only be restricted to this support.

The old source finite correction state has the two exact mean-velocity
constraints
\[
 \int_0^\infty r^2\overline v\,dr=0,\qquad
 \int_0^\infty r\overline\gamma\,dr=0 .
\]
These are source (9.10) in physical variables. The present initial
addition $w_y$ leaves $(\beta,v,\gamma)$ unchanged because its angular
mean is zero. The temporal mean operation has zero auxiliary mean,
and the first two rows of the five-row operation below give zero
increments of these integrals. Thus all specified operations preserve
their initially zero values. Preservation, by itself, does not change
a nonzero entering value to zero.

The already calculated raw tangential residual moments are
\[
 M_2=\partial_zJ_2,\qquad M_1=\partial_zJ_1,\qquad
 J_2=\int_0^\infty r^2C_{z\theta}\,dr,\quad
 J_1=\int_0^\infty rC_{zz}\,dr.
\]
Here the candidate has been added with its zero-angular-mean pressure
$\pi_y$, so the mean pressure is still the old mean pressure.
The source algorithm next reconstructs the mean pressure from the
updated velocity. This operation changes the axial residual moment.

To compute it, first let $\mathcal C_{ab}$ denote the angular average
of the tensor inside \eqref{momop:C}, retaining its auxiliary dependence.
Let $\mathcal D_r$ be the physical radial derivative on the extended
domain, including the auxiliary phase-map derivative. Its auxiliary
average is ordinary $\partial_r$. Since the mean velocity is unchanged
by $w_y$, subtracting the radial conservative equations gives
\begin{equation}\label{momop:delta_g}
 \Delta g_r=-(\mathcal D_r+r^{-1})\mathcal C_{rr}
                   -\partial_z\mathcal C_{zr}
                   +r^{-1}\mathcal C_{\theta\theta}.
\end{equation}
Indeed the radial transport of a divergence-free velocity is
$(\mathcal D_r+r^{-1})u_r^2+
r^{-1}\partial_\theta(u_\theta u_r)+
\partial_z(u_zu_r)-u_\theta^2/r$.
Angular averaging removes the angular derivative and all products
containing exactly one old mean factor and one zero-angular-mean wave.
The source defines $g_r$ as the negative of the mean radial residual
apart from the pressure gradient, which gives exactly the signs in
\eqref{momop:delta_g}. The time and viscous averages of $w_y$
are zero, so no physical viscosity has been dropped in this subtraction.
The formula remains true for any fixed $\nu_{\rm NS}>0$ in the
physical residual.

Set $g=\overline{\Delta g_r}$. Auxiliary derivative integrals vanish,
and hence
\begin{equation}\label{momop:gbar}
 g=-(\partial_r+r^{-1})C_{rr}-\partial_zC_{zr}
                                    +r^{-1}C_{\theta\theta}.
\end{equation}
The increment of the source's pressure defect is therefore the
following explicit new scalar, rather than either raw tangential moment:
\begin{equation}\label{momop:deltaP}
 \Delta P=\int_0^\infty g\,dr
 =\int_0^\infty\frac{C_{\theta\theta}-C_{rr}}r\,dr
                           -\partial_z\int_0^\infty C_{zr}\,dr .
\end{equation}
The equality integrates the derivative $\partial_rC_{rr}$;
both endpoints vanish by compact support. The weighted second
moment of this same source is
\begin{equation}\label{momop:gsecond}
 \int_0^\infty r^2g\,dr
 =\int_0^\infty r(C_{rr}+C_{\theta\theta})\,dr
                   -\partial_z\int_0^\infty r^2C_{zr}\,dr .
\end{equation}
For its first term, integration by parts gives
$-\int r^2\partial_rC_{rr}=2\int rC_{rr}$, after which the
explicit $-C_{rr}/r$ contribution subtracts one copy.
These two identities retain the radial and axial coupling exactly.

\subsection{Compact mean pressure and the corrected axial flux moment}

Retain the source concentration scale and its full axial derivative:
\[
 z=q^D\eta,\qquad 1-t=q(1-\eta^2),\qquad
 A=\tfrac12+h,\quad D=\tfrac12-h,\quad
 q_z=\frac{2\eta q^{1-D}}{1-2h\eta^2}.
\]
Choose the source fixed bump $\widehat\rho$ in its reserved mean
interval $I_m$, with $\int\widehat\rho(x)\,dx=1$, and put
\[
 \rho_{\rm phys}=q^{-1/2}\widehat\rho(r/\sqrt q),\qquad
 c_{\rho,\rm phys}=\frac12\int r^2\rho_{\rm phys}\,dr
             =\frac q2\int x^2\widehat\rho(x)\,dx .
\]
Thus $(c_{\rho,\rm phys})_z=(q_z/q)c_{\rho,\rm phys}$;
it is generally not constant in $z$.

For precision the source's compact physical primitive is as follows.
The fixed slow cutoff $\chi_m$ is zero on the inner collar and one
on the outer collar of the active annulus. For an auxiliary-dependent
scalar $f$ define
\[
 \begin{split}
 I f(r,Y)&=\int_0^r
 f(s,Y+(s^{d_r}-r^{d_r})v_r)\,ds,\\
 J f(r,Y)&=\int_0^\infty
 f(s,Y+(s^{d_r}-r^{d_r})v_r)\,ds,\qquad
 I_cf=If-\chi_mJf.
 \end{split}
\]
Here $d_r=2((1+h)\rho_g-h\kappa_s)$ is the exact source exponent,
$\rho_g=\log(4-\sqrt2)/\log(4+\sqrt2)$,
$v_r=(1,-(\sqrt2-1))$, and $\kappa_s=10^{-5}$.
The derivative $\mathcal D_r$ of the shifted integral differentiates
the integration endpoint while cancelling the derivative of its
$-r^{d_r}v_r$ shift. Consequently
$\mathcal D_rI_cf=f-(\partial_r\chi_m)Jf$.
Also $\overline{Jf}=\int\overline f\,dr$, by translation invariance
of Haar measure. This proves the mean identities used next directly.

The source reconstruction after the actual wave addition changes the
physical mean pressure by
\[
 \Delta p_m=I_c(\Delta g_r-\rho_{\rm phys}\Delta P).
\]
Its source has zero radial auxiliary-mean integral by
\eqref{momop:deltaP}. Therefore
\begin{equation}\label{momop:pressureidentity}
 \partial_r\overline{\Delta p_m}
     =g-\rho_{\rm phys}\Delta P,\qquad
 \int r\,\overline{\Delta p_m}\,dr
     =-\frac12\int r^2g\,dr+c_{\rho,\rm phys}\Delta P .
\end{equation}
The first identity follows from the displayed primitive formula;
its cutoff remainder has zero auxiliary average. The second integrates
$\partial_r(r^2\overline{\Delta p_m})$ between the two zero endpoints.
The full radial residual increment after this pressure update is
$-\rho_{\rm phys}\Delta P-
(\partial_r\chi_m)J(\Delta g_r-\rho_{\rm phys}\Delta P)$.
The second term is retained. The source proves its terminal flatness
under its all-order uniform class bounds; finite smoothness alone is
not that proof.

Use different symbols for the source's two flux defects and the raw
candidate fluxes $J_2,J_1$. Their increments are exactly
\begin{align}
 \Delta J_\theta&=J_2,\nonumber\\
 \Delta J_z&=J_1-\frac12\int r^2g\,dr\nonumber\\
 &=\int r\left(C_{zz}-\frac12C_{rr}
                         -\frac12C_{\theta\theta}\right)\,dr
                    +\frac12\partial_z\int r^2C_{zr}\,dr .
 \label{momop:defectdictionary}
\end{align}
These are the differences of source (8.15): the mean-velocity terms
are unchanged and the full wave covariance changes by $C$.
The axial pressure identity \eqref{momop:pressureidentity} now proves
the exact moments after pressure reconstruction:
\begin{equation}\label{momop:afterpressure}
 M_{2,\rm new}=\partial_z\Delta J_\theta=M_2,\qquad
 M_{1,\rm new}=\partial_z(\Delta J_z+c_{\rho,\rm phys}\Delta P)
  =M_1+\partial_z\left(-\frac12\int r^2g\,dr
                                      +c_{\rho,\rm phys}\Delta P\right).
\end{equation}
These are increment identities. The old source moments and defects
are added to obtain the full updated state. No hypothesis that the
old defects vanish was used: identical old mean-velocity terms
cancel in the subtraction. The product involving $c_{\rho,\rm phys}$
stays inside $\partial_z$, including its nonzero derivative.

\subsection{New quantitative budgets for the three actual inputs}

Retain the previously proved quantities
$\mathcal W_e,\mathcal W_{e,z},\mathcal V_e,\mathcal Z_e$ for
$e=1,2$. They satisfy
\[
 \|w_*\|_{[e]}=\mathcal W_e,\quad
 \|\partial_zw_*\|_{[e]}=\mathcal W_{e,z},\quad
 \|w_y\|_{[e]}\le\mathcal V_e,\quad
 \|\partial_zw_y\|_{[e]}\le\mathcal Z_e.
\]
In particular $\mathcal V_e,\mathcal Z_e$ retain separately the
actual amplitude bounds
$Y_1=(A_0/\mu)\widehat\lambda_3^{-7/8}$ and
$Y_2=2A_0\widehat\lambda_3^{1/8}/(\sigma_2\sqrt{n_i})$
and their physical source-band factors. Define merely for the
following displayed estimates
\[
 B_e=2\mathcal W_e\mathcal V_e+\mathcal V_e^2,\qquad
 D_e=\mathcal W_{e,z}\mathcal V_e+
                    \mathcal W_e\mathcal Z_e+\mathcal V_e\mathcal Z_e.
\]
Choose $r_->0$ enclosing the radial support of the new fields and
their axial derivatives from below. Cauchy--Schwarz on the original
weighted product measures gives
\begin{equation}\label{momop:targetbounds}
 |\Delta P|\le r_-^{-2}B_1+2r_-^{-1}D_1,\qquad
 |\Delta J_\theta|\le B_2,\qquad
 |\Delta J_z|\le B_1+D_2.
\end{equation}
Here is the full estimate. The matrix
$\operatorname{diag}(-1,1,0)$ has operator norm one, so
$|C_{\theta\theta}-C_{rr}|$ is bounded after averaging by
$2|w_*||w_y|+|w_y|^2$. Replacing $r^{-1}$ by
$r_-^{-2}r$ and applying Cauchy--Schwarz gives the first part
$r_-^{-2}B_1$ of the pressure bound. Differentiating $C_{zr}$
produces the six actual cross and self product terms. Their integral
with weight $r^e$ is at most $2D_e$. Replace the unweighted
integral by $r_-^{-1}$ times the $r$-weighted integral to obtain
$2r_-^{-1}D_1$ in \eqref{momop:deltaP}.
The same product estimate without differentiation bounds
$\int r^2C_{z\theta}$ by $B_2$.
Finally $\operatorname{diag}(-1/2,-1/2,1)$ has norm one.
Apply it to the first integral in \eqref{momop:defectdictionary}
to get $B_1$; the last differentiated integral is bounded by
$\tfrac12(2D_2)=D_2$. This proves all three claims without
replacing the two amplitude components by a new unspecified constant.

\subsection{The exact five-row map for these targets}

The source map acts on the \emph{current} physical triple
$(P,J_\theta,J_z)$, after adding \eqref{momop:deltaP} and
\eqref{momop:defectdictionary} to the old source triple and
reconstructing pressure. It does not take
$(M_2,M_1)$ as its input. The five equations, in physical variables, are
\begin{align}
 \int r^2\Delta v\,dr&=0,&
 \int r\gamma_d\,dr&=0,\nonumber\\
 \int\frac{2V}{r}\Delta v\,dr&=-P,\nonumber\\
 \int r^2(G\Delta v+V\gamma_d)\,dr&=-J_\theta,\nonumber\\
 \int(2rG\gamma_d-rV\Delta v)\,dr&=-J_z .
 \label{momop:fivephysical}
\end{align}
In the last row there is exactly one factor $r$ multiplying $V\Delta v$.
The pressure-source term contributes
$-\tfrac12\int r^2(2V/r)\Delta v\,dr=-\int rV\Delta v\,dr$,
which proves that factor rather than treating the fifth row as an
independent constraint with an arbitrary weight.

Retain the source's reserved patch $x=r/\sqrt q\in I_m$. Its actual
summed base satisfies source (5.44) and (8.24):
\[
\begin{gathered}
 V_{\rm phys}=q^{-A}a(\eta)x^{-1-2\lambda},\qquad
 G_{\rm phys}=0,\\
 a(\eta)=2^{1/2+\lambda}c_{\rm patch}(1+\eta^2)^{-1},
 \quad \lambda>0,\quad |a(\eta)|\ge a_0>0.
\end{gathered}
\]
The source fixed $\lambda$ here is its reserved-profile exponent,
not a frequency of the coupled problem. In a fixed $Q$ chart,
the same physical base is
\[
 V_*=(Q/q)^A a(\eta)\bigl(\sqrt{Q/q}\,R\bigr)^{-1-2\lambda},
 \qquad G_*=0.
\]
The factor $(Q/q)^A$ has been retained.

Choose a nonnegative nonzero bump $\varphi_0\in C_c^\infty(I_m)$
supported near an interior point and a fixed spacing $d_b>0$ so small
that the three profiles
\[
 \varphi_j(x)=e^{-jd_b}\varphi_0(e^{-jd_b}x),\qquad j=0,1,2,
\]
have disjoint supports in the patch. This is possible because the patch
contains an open interval about that point: choose three distinct
scaled centers there, then decrease the support radius below half the
minimum separation. The symbol $d_b$ distinguishes this source bump
spacing from the coupled physical diffusivity $d$.
For every real $p$, set
\[
 \mu_p=\int x^p\varphi_0(x)\,dx>0,\qquad
 \int x^p\varphi_j(x)\,dx=\mu_pe^{jd_bp}.
\]
The equality is the substitution $x=e^{jd_b}\xi$ and retains the
prefactor in $\varphi_j$.

Define the two explicit moment matrices
\[
 A_\theta=
 \begin{pmatrix}
 \mu_2&\mu_2e^{2d_b}&\mu_2e^{4d_b}\\
 2a\mu_{-2-2\lambda}&
 2a\mu_{-2-2\lambda}e^{d_b(-2-2\lambda)}&
 2a\mu_{-2-2\lambda}e^{2d_b(-2-2\lambda)}\\
 -a\mu_{-2\lambda}&-a\mu_{-2\lambda}e^{-2d_b\lambda}&
                              -a\mu_{-2\lambda}e^{-4d_b\lambda}
 \end{pmatrix},
\]
\[
 A_z=
 \begin{pmatrix}
 \mu_1&\mu_1e^{d_b}\\
 a\mu_{1-2\lambda}&a\mu_{1-2\lambda}e^{d_b(1-2\lambda)}
 \end{pmatrix},\qquad a=a(\eta).
\]
Both inverses exist at every source parameter value. Indeed the
angular matrix is a diagonal row factor times the matrix with rows
$(1,t_p,t_p^2)$ for $p=2,-2-2\lambda,-2\lambda$ and
$t_p=e^{d_bp}$. Subtracting the first row from the other two
and factoring their differences gives
\[
 \det[1,t_p,t_p^2]_{p=p_1,p_2,p_3}
 =(t_{p_2}-t_{p_1})(t_{p_3}-t_{p_1})(t_{p_3}-t_{p_2}).
\]
All three powers are distinct for $\lambda>0$. The diagonal
row factor has determinant
\[
-2a^2\mu_2\mu_{-2-2\lambda}\mu_{-2\lambda}\ne0.
\]
Likewise
\[
\det A_z=a\mu_1\mu_{1-2\lambda}
(e^{d_b(1-2\lambda)}-e^{d_b})\ne0.
\]
This also gives an explicit finite inverse norm bound
$\|A^{-1}\|\le\|\operatorname{adj}A\|/|\det A|$
for either matrix. No uniform limit as $\lambda\downarrow0$ is taken.

Use the exact row normalizations
\[
 P_q=q^{2A}P,\qquad
 J_{\theta,q}=q^{2A-3/2}J_\theta,\qquad
 J_{z,q}=q^{2A-1}J_z,
\]
and define the five coefficients by
\begin{equation}\label{momop:coefficients}
 {\bf u}=A_\theta^{-1}(0,-P_q,-J_{z,q})^T,\qquad
 {\bf s}=A_z^{-1}(0,-J_{\theta,q})^T .
\end{equation}
Then the actual physical velocity increments are
\begin{equation}\label{momop:fiveincrement}
 \Delta v=q^{-A}\sum_{j=0}^2u_j\varphi_j(r/\sqrt q),
 \qquad
 \gamma_d=q^{-A}\sum_{j=0}^1s_j\varphi_j(r/\sqrt q).
\end{equation}
The first and second velocity moments have factors $q^{3/2-A}$
and $q^{1-A}$, respectively. The three inhomogeneous rows have
factors $q^{-2A}$, $q^{3/2-2A}$, and $q^{1-2A}$, respectively.
Substituting the original base into \eqref{momop:fivephysical}
and changing variables $r=\sqrt q\,x$ therefore gives exactly
the two linear systems \eqref{momop:coefficients}. This proves all
five equations, uniqueness in the selected five-dimensional span,
and smooth dependence on the actual target values. For instance
$\partial_\xi A^{-1}=-A^{-1}(\partial_\xi A)A^{-1}$ follows by
differentiating $AA^{-1}=I$; repeated product rules give every
coefficient derivative with the actual $q,\eta$ factors retained.

The desired axial increment has zero weighted integral and is
auxiliary independent. Its exact compact azimuthal potential and
radial increment are
\[
 \Psi(r,z,t)=\frac1r\int_0^r s\gamma_d(s,z,t)\,ds,\qquad
 \Delta\beta=-\partial_z\Psi,\qquad
 \Delta\gamma=(\partial_r+r^{-1})\Psi=\gamma_d.
\]
The zero moment in \eqref{momop:fivephysical} makes $\Psi$ vanish
beyond its outer support; its lower-end definition makes it vanish
below the inner support. Differentiation proves
$(\partial_r+r^{-1})\Delta\beta+\partial_z\Delta\gamma=0$.
The support includes the intervals between bump supports, which still
lie within the reserved patch.
To display the axial scaling cost explicitly, put
$\Psi_j(x)=x^{-1}\int_0^x \xi\varphi_j(\xi)\,d\xi$. Then
\begin{align}
 \Psi&=q^{1/2-A}\sum_{j=0}^1s_j\Psi_j(x),\nonumber\\
 \Delta\beta
 &=-q^{1/2-A}\sum_{j=0}^1\left[
    (\partial_zs_j)\Psi_j+
    \frac{q_z}{q}s_j\left((1/2-A)\Psi_j
                                      -\frac x2\Psi_j'\right)\right].
 \label{momop:inducedradial}
\end{align}
The two individual $\Psi_j$ can have $x^{-1}$ tails. The exact
linear relation $\sum_js_j\mu_1e^{jd_b}=0$ cancels those tails
in their sum and, after differentiation, in \eqref{momop:inducedradial}.
No individual compactness was assumed.

\subsection{What the five-row operation actually leaves}

Let $(\beta,v,\gamma)$ be the current mean correction, after whatever
preceding wave and mean updates have actually been performed. Keep the
base $(b,V,G)$ and the full current wave $w$ fixed during this
five-row step. Use the three full current defects as its targets.
Write $\mathcal D_r$ for the physical extended radial derivative and
$\mathcal D_t$ for the physical extended time derivative. With
$\Delta_0=\mathcal D_r^2+r^{-1}\mathcal D_r+\partial_z^2$, direct
subtraction of the full radial source gives
\begin{align}
 R_g:={}&g_{r,\rm new}-g_r-\frac{2V}{r}\Delta v\nonumber\\
 ={}&-\mathcal D_t\Delta\beta
 -(\mathcal D_r+r^{-1})
             \bigl(2b\Delta\beta+2\beta\Delta\beta+(\Delta\beta)^2\bigr)
 \nonumber\\
 &-\partial_z\bigl(b\Delta\gamma+G\Delta\beta+
           \beta\Delta\gamma+\gamma\Delta\beta+\Delta\beta\Delta\gamma\bigr)
 \nonumber\\
 &+\frac{2v\Delta v+(\Delta v)^2}{r}
               +\nu_{\rm NS}(\Delta_0-r^{-2})\Delta\beta .
 \label{momop:Rg}
\end{align}
Here all increments are auxiliary independent, so their
$\mathcal D_t$ is ordinary $\partial_t$. The products with the
old means still retain their auxiliary dependence. Equation
\eqref{momop:Rg} keeps the actual physical viscosity; the source
unit-viscosity formula (8.26) is its case $\nu_{\rm NS}=1$.

The full updated defects, with pressure recomputed from the new
radial source, are exactly
\begin{align}
 P_{\rm new}&=\int\overline{R_g}\,dr,\nonumber\\
 (J_\theta)_{\rm new}
    &=\int r^2\overline{\gamma\Delta v+v\Delta\gamma+
                                      \Delta\gamma\Delta v}\,dr,\nonumber\\
 (J_z)_{\rm new}
    &=\int r\overline{2\gamma\Delta\gamma+(\Delta\gamma)^2}\,dr
                      -\frac12\int r^2\overline{R_g}\,dr .
 \label{momop:recomputed}
\end{align}
To verify these formulas, expand every quadratic product in the
definitions of the current and new defects. The radial source
increment is $(2V/r)\Delta v+R_g$. The third row of
\eqref{momop:fivephysical} cancels the old $P$ with the integral
of the displayed linear term. The azimuthal flux increment contains
the base-linear terms $G\Delta v+V\Delta\gamma$, cancelled by
the fourth row because $\Delta\gamma=\gamma_d$. The axial defect
contains $2G\Delta\gamma$ and the radial-source contribution
$-\tfrac12r(2V\Delta v)$ after integration. These are precisely
the fifth row. All remaining products are the terms displayed
in \eqref{momop:recomputed}; the full wave covariance is unchanged
during this mean step. This proves the formulas without assuming
their right-hand sides vanish.

The two velocity moments remain exactly zero. The new residual
moments therefore have the source form
\[
 \int r^2\overline{E_{\theta,\rm new}}\,dr
       =\partial_z(J_\theta)_{\rm new},\qquad
 \int r\overline{E_{z,\rm new}}\,dr
       =\partial_z\bigl((J_z)_{\rm new}
                           +c_{\rho,\rm phys}P_{\rm new}\bigr).
\]
The pressure and $c_{\rho,\rm phys}$ are the actual reconstructed
objects. Thus the five equations cancel specified linear defect
contributions and expose the remaining nonlinear and differential
terms. Equations \eqref{momop:Rg}--\eqref{momop:recomputed} supply
the next actual inputs rather than a claim of complete cancellation.

\subsection{Exact source operation dictionary and the bounds it requires}

For the source's estimates, $Q=2^{-\ell}$, $\varepsilon=Q^h$,
$S_*=\ell^2$, $R=r/\sqrt Q$, $Z=z/Q^D$, $T=(1-t)/Q$.
Use a coarsest applicable common auxiliary torus $H$ as in source
Lemma 6.2; different representatives agree by its integer covering
maps. The original radial edge weights are
\[
\begin{gathered}
 \zeta(X)=\exp\left(-\frac{a_a}{\log^2(X/X_a)}
                       -\frac{a_b}{\log^2(X_b/X)}\right),\\
 \delta(X)=\min\{1,\log(X/X_a),\log(X_b/X)\},
 \quad X=\frac{QR^2}{2q},
\end{gathered}
\]
with the source fixed $a_a,a_b>0$. In the following bounds the band,
label, representative and harmonic are held fixed when differentiating.
The constants may depend on the finite correction stage and derivative
multi-index, but are uniform over bands, labels, points and rectangle
copies.

A mean coefficient $f\in M^\alpha$ has no angular dependence, descends
to the common torus, extends smoothly by zero outside the active shell,
and satisfies for each multi-index $I$ in $(R,Z,T,H)$
\[
 |\partial^If|\le C_I\varepsilon^\alpha S_*^{b_I}
                                        \zeta\,\delta^{-d_I}.
\]
It can have nonzero auxiliary modes and has no rectangle or pulse
envelope restriction. A moment coefficient $F(Z,T)\in S^\alpha$
satisfies $|\partial_{Z,T}^IF|\le C_I\varepsilon^\alpha S_*^{b_I}$.
A labelled nonzero harmonic amplitude $a_{\gamma,m}\in W^\alpha$
has no angular dependence, the source slow, transverse and shell
support, smooth zero extension across those boundaries, compatible
lift representatives, and
\[
 |\partial^Ia_{\gamma,m}|\le
 C_I\varepsilon^\alpha S_*^{b_I}\sqrt\zeta\,\delta^{-d_I}
                          {\cal P}_\gamma(v),\qquad 0\le v\le L_s.
\]
The prescribed source phase has $k_\gamma p_\gamma\in
\mathbb Z\setminus\{0\}$ and $k_\gamma=\lceil\varepsilon^{-1/2}\rceil$.
The harmonic set is finite at a fixed stage and independent of band.
Equal label--harmonic terms are grouped before testing the bounds.
The pulse envelope is the original positive exponential in (7.12),
with its proved Gaussian estimates. Actual cut-off waves additionally
have support where the source pulse cutoff $\psi$ is supported and
smooth zero extension in $v$. No divisibility by $\psi$ is required.

The source's anisotropic derivative mappings, with the phase factored
out, are
\[
\begin{gathered}
 \partial_{R,Z,T}:C^\alpha\to C^\alpha,\quad
 D_r:C^\alpha\to C^{\alpha-\kappa_s},\quad
 D_z=\varepsilon\partial_Z:C^\alpha\to C^{\alpha+1},\\
 c_{i_0}N_{i_0}:C^\alpha\to C^\alpha,\quad
 -\varepsilon\partial_T:C^\alpha\to C^{\alpha+1},
 \qquad C=M,W.
\end{gathered}
\]
Products obey $M^\alpha M^\beta\subset M^{\alpha+\beta}$ and
$M^\alpha W^\beta\subset W^{\alpha+\beta}$. Same-label wave
products have this summed order, in $M$ for zero output harmonic
and $W$ otherwise; distinct supported labels have zero products.
The zero-harmonic global $M$ statement here is for the actual temporally
cut-off waves. Before that cutoff, the source gives the mean bounds on
the labelled rectangle, not a global $M$ assertion.
These are source Proposition 6.6, not deductions from the candidate's
finite norm bounds.

Every radial target must first carry the correct measure and weight.
For $f_*=Q^{a_f}f_{\rm phys}$,
\[
 \int R^e\overline{f_*}\,dR
       =Q^{a_f-(e+1)/2}\int r^e\overline{f_{\rm phys}}\,dr .
\]
For the present inputs this reads
\begin{align*}
 P_*&=Q^{2A}P,&
 (J_\theta)_*&=Q^{2A-3/2}J_\theta,&
 (J_z)_*&=Q^{2A-1}J_z,&
 (c_\rho)_*&=Q^{-1}c_{\rho,\rm phys},\\
 (M_2)_*&=Q^{2A-1}M_2,&
 (M_1)_*&=Q^{2A-1/2}M_1 .
\end{align*}
In particular \eqref{momop:afterpressure} becomes exactly
$\varepsilon\partial_Z\Delta J_{\theta,*}$ and
$\varepsilon\partial_Z(\Delta J_{z,*}+c_{\rho,*}\Delta P_*)$.
The factors match because $h+D=1/2$. This is the extra axial
factor used in source (9.12); using the raw $J_1$ after pressure
reconstruction would omit part of its actual input.

For clarity, the following are the four ordered operations of source
Proposition 9.6, with the full current velocity and pressure
reconstructed after each operation.

\begin{enumerate}
\item For each supported nonzero harmonic coefficient $f_{\gamma,m}$
of the full current residual, solve the original transverse pulse
equation with zero entrance data and source $-f_{\gamma,m}$.
In frame coordinates its coefficient matrix is
$H_\gamma-m^2\varepsilon k_\gamma^2|n_\gamma|^2I$.
Multiply the solved potential and pressure by $\psi$ and take the
full curl. Source Proposition 7.2 maps $W^\alpha$ sources to
$W^\alpha$ velocity amplitudes and $W^{\alpha+1/2}$ pressure.
Proposition 9.1 assigns the remaining supported linear residual
order $\alpha+1/2-3\kappa_s$ and retains its Gaussian cutoff tails
additively. The actual candidate's full residual formula determines
the input coefficients. Following this wave change, replace $w_y$
by the actual total new wave increment in \eqref{momop:C}, recompute
$g_r$ and reconstruct pressure. Its old-wave cross terms must change
with the updated wave.

\item Let $F_2,F_1$ be the auxiliary-averaged physical tangential
residuals of that updated state. Their two moments are the updated
versions of \eqref{momop:afterpressure}, including the old defects.
For each $e=1,2$ retain the fixed normalized physical bump
$b_e=q^{-(e+1)/2}\widehat b_e(r/\sqrt q)$ with
$\int x^e\widehat b_e(x)\,dx=1$. Form
\[
 \mathfrak M_e=\int r^eF_e\,dr,\qquad
 \sigma_e=-r^{-e}\int_0^r s^e
                         (F_e(s)-b_e(s)\mathfrak M_e)\,ds .
\]
The zero total weighted integral proves compact support and
$(\partial_r+e/r)\sigma_e=-F_e+b_e\mathfrak M_e$.
Use $\Sigma=Q^{2A}(\sigma_2,\sigma_1)$ as the input to the
source's signed amplitude derivative (7.31), with its fixed
positive primary amplitudes. Source Proposition 7.6 maps an
auxiliary-independent $M^\alpha$ stress to $W^{\alpha-1/2}$.
Take its actual curl, retain the old-wave cross and quadratic
correction tensor (9.13), and reconstruct pressure. The source
does not set $b_e\mathfrak M_e$ to zero.

If $\mathfrak M_e=\partial_z\mathfrak J_e$, its bump term satisfies
the exact derivative identity
\[
 b_e\partial_z\mathfrak J_e
 =\partial_z(b_e\mathfrak J_e)-(\partial_zb_e)\mathfrak J_e,\qquad
 \partial_zb_e=-\frac{q_z}{2q}
                    \bigl((e+1)b_e+r\partial_rb_e\bigr).
\]
It follows by differentiating the original $q$-scaled bump.
Thus trading the moment bump for an axial tensor flux leaves the
displayed radial-profile derivative term. For the axial equation,
$\mathfrak J_1=J_z+c_{\rho,\rm phys}P$, including both product
derivatives in \eqref{momop:afterpressure}.

\item For the zero-auxiliary-mean part of the current angularly
invariant residual, first convert the physical residual to the fixed
source chart. In this step a star denotes that chart representative:
\[
\begin{gathered}
 E_*=Q^{2A+1/2}E_{\rm phys},\qquad
 E_{a,*}^\circ=E_{a,*}-\langle E_{a,*}\rangle_Y,\\
 (\Delta\beta_*,\Delta v_*,\Delta\gamma_*)
       =Q^A(\Delta\beta_{\rm phys},\Delta v_{\rm phys},
                                      \Delta\gamma_{\rm phys}).
\end{gathered}
\]
The source temporal update in these exact units is
\[
 \Delta v_*=-c_{i_0}^{-1}N_{i_0}^{-1}E_{\theta,*}^\circ,\qquad
 \gamma_{d,*}=-c_{i_0}^{-1}N_{i_0}^{-1}E_{z,*}^\circ,
 \quad
 N_{i_0}^{-1}F=
 \sum_{k\ne0}\frac{\widehat F(k)}{2\pi i v_t\cdot k}e^{2\pi ik\cdot H}.
\]
Here $v_t=(\sqrt2-1,1)$ and
$c_{i_0}=T_g^{i_0}Q^{1+h}$, with $c_{i_0}^{-1}\le CS_*$.
The source Diophantine bound gives
$\|N^{-1}F\|_{C^m}\le C_m\|F\|_{C^{m+4}}$.
Let $T_{1,*}$ and $A_{1,*}$ be the source's radial potential and
its cutoff remainder in this chart. Realize $\gamma_{d,*}$ by
$T_{1,*}\gamma_{d,*}$:
$\Delta\beta_*=-\varepsilon\partial_ZT_{1,*}\gamma_{d,*}$ and
$\Delta\gamma_*=\gamma_{d,*}-A_{1,*}\gamma_{d,*}$.
The actual axial remainder $-A_{1,*}\gamma_{d,*}$ stays in the
new velocity. The exact normalized cancellations are
\[
\begin{split}
 c_{i_0}N_{i_0}\Delta v_*&=-E_{\theta,*}^\circ,\\
 c_{i_0}N_{i_0}\Delta\gamma_*&=-E_{z,*}^\circ
              -c_{i_0}N_{i_0}A_{1,*}\gamma_{d,*}.
\end{split}
\]
Convert every resulting velocity component back by $Q^{-A}$ before
reconstructing the physical pressure or taking physical products and
defects. In particular its physical angular update is exactly
\[
 \Delta v_{\rm phys}
  =-Q^{-A}c_{i_0}^{-1}Q^{2A+1/2}
                     N_{i_0}^{-1}E_{\theta,\rm phys}^\circ
  =-T_g^{-i_0}N_{i_0}^{-1}E_{\theta,\rm phys}^\circ .
\]
The last equality uses $A+1/2=1+h$ and the original
$c_{i_0}=T_g^{i_0}Q^{1+h}$. It is also
$-N_{\rm abs}^{-1}E_{\theta,\rm phys}^\circ$, since
$N_{\rm abs}=T_g^{i_0}N_{i_0}$ on the representative.
This proves the conversion without applying a normalized inverse
directly to a physical residual.

\item Apply \eqref{momop:coefficients} to the three defects of this
updated state, use \eqref{momop:fiveincrement} and its actual
radial component \eqref{momop:inducedradial}, and reconstruct
pressure again. The exact remaining defects are
\eqref{momop:Rg}--\eqref{momop:recomputed}.
Source Lemma 8.7 maps $S^\alpha$ targets to $M^\alpha$
tangential increments and $M^{\alpha+1}$ radial increments.
Source Lemma 8.8 gives the remaining defect class
$S^{\alpha+0.9-2\kappa_s}$ using the actual bounds
$v,\gamma\in M^{0.9}$, $\beta\in M^{1.9}$, $\alpha\ge0.9$,
and its reserved-patch base derivative bounds.
\end{enumerate}

The source cycle assumes its full finite correction-state bounds
$w\in W^{1/2}$, $w-w_0^{\tan}\in W^{0.68}$,
$v,\gamma,p_m\in M^{0.9}$, $\beta\in M^{1.9}$,
with supported nonzero residual order
$B_j=1/2+\sigma_j$, mean and defect order $C_j^*=1+\sigma_j$,
and $\sigma_j=1/5+j/10$. Its ordered estimate proves gains
$B_{j+1}-B_j=C_{j+1}^*-C_j^*=1/10$.
These are requirements and conclusions for the source's construction.
The explicit candidate budgets and equations above do not by
themselves identify those exponents.

There is a precise finite-domain statement available now. Write the
full current residual as $R_{\rm old}+\Delta R$, retaining the old
source state and every term of the computed increment. The added
candidate uses finitely many labels, with its cutoffs supported
strictly inside the source rectangle and radial annulus. On those
compact coefficient supports, $\varepsilon>0$, $\zeta>0$,
$\delta>0$ and $\mathcal P_\gamma(v)>0$ have positive minima
for each fixed label. By local finiteness of the retained source-label
family, only finitely many old labels intersect this fixed compact
perturbation support. Every differentiated coefficient of $\Delta R$
is finite by its explicit smooth shapes, including the retained old-state
cross terms there. For any fixed
$\alpha$ and derivative $I$, the maximum of its derivative divided by
$\varepsilon^\alpha\sqrt\zeta\mathcal P_\gamma$ on these finitely
many supports is therefore a concrete finite constant. Extend the
finite \emph{increment} family by zero on unused labels. This proves its
finite coefficient inequalities without assigning an asymptotic gain
to the amplitude. It does not set any old source coefficient to zero:
$R_{\rm old}$ retains its actual locally finite family and its attributed
source-class bounds. The finite-sum closure in source Proposition 6.6
combines those bounds with the proved increment bounds for each full
current input. The same compact-support argument applies to the added
mean sources and the reserved bump operators. Thus the operations
above are defined on the actual old-plus-increment inputs.
Their constants can contain the inverse of a very small Gaussian
envelope or source scale; no uniform control for an increasing family
of coupled stages or bands follows from this compactness argument.
The original component budgets and the full nonlinear remainders
remain the quantities to propagate when investigating such a family.


\section{Compact completion of the full averaged residual}
\label{mcstress:section}

This section computes a compact symmetric stress and a compact pressure
for the complete averaged residual of the actual finite curl candidate.
It retains the original radial bumps depending on $q(z,t)$, both axial
fluxes, the radial force, and the full cylindrical connection terms.
The result is a stress--pressure representation of a computed force;
it does not by itself modify a velocity or identify the constructed
stress with a covariance of additional waves.

\subsection{The actual extended fields and all three averaged components}

Use right-handed cylindrical coordinates
$x=(r\cos\theta,r\sin\theta,z)$, with oriented orthonormal basis
$(e_r,e_\theta,e_z)$ and $r>0$. The retained finite source state is
$(u_*,p_*)$; the actual added field and its selected pressure are
$(v,\pi)=(w_y,\pi_y)$. On a finite collection of source labels the
added potential has the specified form
\begin{equation}\label{mcstress:actual_potential}
\begin{gathered}
 \mathcal A_y=\sum_j\operatorname{Re}
   \left[Q_j^{1/2-A}
     \frac{i\,n_j\times t_j}{k_jm_j|n_j|^2}
                      e^{ik_jm_j\Phi_j}\right],\qquad
 v=\operatorname{curl}_{\mathscr D}\mathcal A_y,\\
 t_j=\chi_jB_jy_3(t_{c,j}),\qquad
 y_3=\binom{\Theta_3}{\Omega_3},\qquad
 t_{c,j}=t_b+\alpha_jv_j,\qquad
 n_j^TB_j=0,\qquad k_jm_jp_j\in\mathbb Z\setminus\{0\}.
\end{gathered}
\end{equation}
Here $y_3$ is the actual signed third-growth trajectory, the cutoffs
and source frame are the retained functions, and $Q_j$ is constant
on each source chart. The chosen pressure is the real harmonic sum
with chart coefficient
\[
 Q_j^{-2A}
 \frac{i\{n_j^T\mathsf K_jt_j-(n'_j)^Tt_j\}}
      {k_jm_j|n_j|^2}e^{ik_jm_j\Phi_j}.
\]
All coefficient components are independent of $\theta$; their angular
dependence is the displayed nonzero integer harmonic. Their radial
and axial supports lie in a fixed compact set with $r$ bounded away
from zero on the finite time interval under consideration. Smooth
zero extension at each source cutoff is retained.

The source absolute auxiliary coordinate is $Y\in\mathbb T^2$ and
the physical evaluation is
\begin{equation}\label{mcstress:extended_derivatives}
\begin{gathered}
 Y(r,t)=v_r r^{d_r}+v_t t\pmod{\mathbb Z^2},\qquad
 v_r=(1,-b_g),\quad v_t=(b_g,1),\quad b_g=\sqrt2-1,\\
 \rho_g=\frac{\log(4-\sqrt2)}{\log(4+\sqrt2)},\quad
 \kappa_s=10^{-5},\quad d_r=2((1+h)\rho_g-h\kappa_s),\\
 \mathscr D_r=\partial_r+d_r r^{d_r-1}v_r\cdot\partial_Y,
 \qquad \mathscr D_t=\partial_t+v_t\cdot\partial_Y,
 \qquad \mathscr D_z=\partial_z.
\end{gathered}
\end{equation}
The angular derivative is $\partial_\theta$.
The exponent $h$ here is the unchanged source chart exponent. These
lifted derivatives commute, obey the product rule, and satisfy
$\mathscr D_r r=1$. Equivalently in Cartesian coordinates they are
$\partial_{x_i}+\sum_\alpha(\partial_{x_i}Y_\alpha)\partial_{Y_\alpha}$;
their commutators vanish because mixed derivatives of $Y$ commute.
Thus evaluation at the physical map preserves every differentiated
identity below.

The actual finite source has its retained representation
$u_* =\operatorname{curl}_{\mathscr D}\mathcal A_*+B_*e_\theta$,
with $\partial_\theta B_*=0$. Commutation of the lifted derivatives
proves $\operatorname{div}_{\mathscr D}\operatorname{curl}_{\mathscr D}=0$.
Also $\operatorname{div}_{\mathscr D}(B_*e_\theta)=0$.
Consequently both $u_*$ and $v$ are divergence free at every independent
$Y$, not only after physical evaluation.

Let the normalized averages act on cylindrical components before
evaluation:
\[
 \langle f\rangle_\theta=\frac1{2\pi}\int_0^{2\pi}f\,d\theta,
 \quad \langle f\rangle_Y=\int_{\mathbb T^2}f\,dY,
 \quad \langle f\rangle=\langle\langle f\rangle_\theta\rangle_Y,
 \quad\int_{\mathbb T^2}1\,dY=1.
\]
Every cylindrical curl operation retains the nonzero angular harmonic
in \eqref{mcstress:actual_potential}. For example its radial and axial
components are $r^{-1}\partial_\theta\mathcal A_z-\partial_z\mathcal A_\theta$
and $(\mathscr D_r+r^{-1})\mathcal A_\theta-r^{-1}\partial_\theta\mathcal A_r$.
Integrating the integer harmonic proves
\begin{equation}\label{mcstress:zero_mean}
 \langle v\rangle_\theta=0,\qquad
 \langle\pi\rangle_\theta=0
 \quad\hbox{at every }(r,z,t,Y).
\end{equation}
This statement includes the complete coefficient curl and every cutoff
derivative, rather than only the principal amplitude.

For a smooth periodic coefficient $f$, integration of its auxiliary
derivatives gives
\begin{equation}\label{mcstress:average_derivatives}
 \langle\mathscr D_r f\rangle=\partial_r\langle f\rangle,
 \quad \langle\mathscr D_t f\rangle=\partial_t\langle f\rangle,
 \quad \langle\partial_z f\rangle=\partial_z\langle f\rangle,
 \quad \langle\partial_\theta f\rangle=0.
\end{equation}
Applying the first identity twice also treats $\mathscr D_r^2$;
the derivative of its variable coefficient has zero auxiliary average
because that coefficient is independent of $Y$.

At arbitrary fixed physical viscosity $\nu_{\rm NS}>0$, the exact
residual increment is
\begin{equation}\label{mcstress:residual_increment}
 \Delta R=\mathscr D_t v+(u_*\cdot\nabla_{\mathscr D})v
 +(v\cdot\nabla_{\mathscr D})u_*
 +(v\cdot\nabla_{\mathscr D})v
 -\nu_{\rm NS}\Delta_{\mathscr D}v+\nabla_{\mathscr D}\pi.
\end{equation}
This is the residual of the actual modified state minus the source
state's residual at the same viscosity. Put
\begin{equation}\label{mcstress:actual_tensor}
 \bar u_* =\langle u_*\rangle_\theta,\qquad
 w_*=u_*-\bar u_*,\qquad
 S=\langle w_*\otimes v+v\otimes w_*+v\otimes v\rangle.
\end{equation}
It is a smooth real symmetric axisymmetric tensor compactly supported
in the annular and axial set at issue. Each summand contains a factor
of the actual compact field $v$.

\begin{theorem}\label{mcstress:full_mean}
The complete averaged residual is the full Cartesian divergence of
the actual tensor $S$:
\begin{equation}\label{mcstress:Ffull}
 \langle\Delta R\rangle
 =F_r e_r+F_2e_\theta+F_1e_z
 =\operatorname{div}S,
\end{equation}
where
\begin{equation}\label{mcstress:all_components}
\begin{aligned}
 F_r&=(\partial_r+r^{-1})S_{rr}-r^{-1}S_{\theta\theta}
                                      +\partial_zS_{rz},\\
 F_2&=(\partial_r+2r^{-1})S_{r\theta}+\partial_zS_{z\theta},\\
 F_1&=(\partial_r+r^{-1})S_{rz}+\partial_zS_{zz}.
\end{aligned}
\end{equation}
\end{theorem}
\begin{proof}
For divergence-free fields, direct Cartesian component expansion
rewrites the three advection terms of \eqref{mcstress:residual_increment}
as the lifted divergence of
$u_*\otimes v+v\otimes u_*+v\otimes v$.
The mean-flow cross tensor has zero angular average by
\eqref{mcstress:zero_mean}, even when $\bar u_*$ depends on $Y$.
Thus its averaged divergence is zero by
\eqref{mcstress:average_derivatives} and the cylindrical divergence
formula.

To check the linear radial component explicitly, let
\[
 \mathscr L_0=\mathscr D_r^2+r^{-1}\mathscr D_r
             +r^{-2}\partial_\theta^2+\partial_z^2.
\]
The complete cylindrical vector Laplacian is
\[
 \begin{aligned}
 (\Delta_{\mathscr D}v)_r&=(\mathscr L_0-r^{-2})v_r
                       -2r^{-2}\partial_\theta v_\theta,\\
 (\Delta_{\mathscr D}v)_\theta&=(\mathscr L_0-r^{-2})v_\theta
                       +2r^{-2}\partial_\theta v_r,\qquad
 (\Delta_{\mathscr D}v)_z=\mathscr L_0v_z.
 \end{aligned}
\]
Each average is zero by \eqref{mcstress:zero_mean} and
\eqref{mcstress:average_derivatives}. The radial pressure average is
$\langle\mathscr D_r\pi\rangle=\partial_r\langle\pi\rangle=0$;
the angular and axial components vanish as well. The averaged time
derivative is zero. These verifications retain the radial connection
term and physical viscosity before taking the average.

Finally, for a symmetric tensor independent of $\theta$ in cylindrical
components, differentiation of
$e_r=(\cos\theta,\sin\theta,0)$ and
$e_\theta=(-\sin\theta,\cos\theta,0)$ gives exactly the three
expressions in \eqref{mcstress:all_components}. One direct verification
is to write its Cartesian matrix as $O_\theta S O_\theta^T$, where
$O_\theta=[e_r,e_\theta,e_z]$, and use
$\partial_x=\cos\theta\partial_r-r^{-1}\sin\theta\partial_\theta$
and $\partial_y=\sin\theta\partial_r+r^{-1}\cos\theta\partial_\theta$.
The derivatives of $O_\theta$ give $-S_{\theta\theta}/r$ in the radial
row, $2S_{r\theta}/r$ in the angular row, and $S_{rz}/r$ in the axial
row. Averaging the lifted version replaces $\mathscr D_r$ by
$\partial_r$. This proves the full Cartesian identity.
\end{proof}

\subsection{The compact fluxes and the original varying radial bumps}

For $e=1,2$, write $\tau_1=z$, $\tau_2=\theta$, and set
\begin{equation}\label{mcstress:fluxes}
 M_e(z,t)=\int_0^\infty r^eF_e(r,z,t)\,dr,
 \qquad
 J_e(z,t)=\int_0^\infty r^eS_{z\tau_e}(r,z,t)\,dr.
\end{equation}
Multiplication of the corresponding row of
\eqref{mcstress:all_components} by $r^e$ makes its radial term
$\partial_r(r^e S_{r\tau_e})$. The radial boundary values vanish
by compact annular support. Therefore
\begin{equation}\label{mcstress:flux_identity}
 M_e=\partial_zJ_e,\qquad
 \int_{\mathbb R}M_e(z,t)\,dz=0\quad(e=1,2).
\end{equation}
Both $J_e$ are smooth and compactly supported in $z$. No claim that
either $M_e$ vanishes pointwise is used.

Retain the original $q$-dependent bumps, for $q(z,t)>0$ on the compact
source patch:
\begin{equation}\label{mcstress:bumps}
 b_e(r,z,t)=q(z,t)^{-(e+1)/2}
             \widehat b_e(r/\sqrt{q(z,t)}),\qquad
 \widehat b_e\in C_c^\infty((0,\infty)),\qquad
 \int_0^\infty x^e\widehat b_e(x)\,dx=1.
\end{equation}
The fixed profiles may be taken to be the already selected ones.
The change of variable $r=\sqrt q\,x$ proves
$\int r^e b_e\,dr=1$ for every $(z,t)$. Choose a fixed annulus
containing the supports of $S$ and of the bumps on the relevant compact
$(z,t)$ set. This is possible because $q$ has positive lower and finite
upper bounds there. Outside the compact axial set where a flux or
force is nonzero all the formulas below extend by zero.

The exact derivatives and cumulative bump are
\begin{equation}\label{mcstress:bump_derivatives}
\begin{gathered}
 \partial_zb_e=-\frac{q_z}{2q}
                  \bigl((e+1)b_e+r\partial_rb_e\bigr),\\
 \beta_e(r,z,t)=\int_0^r s^eb_e(s,z,t)\,ds
              =\int_0^{r/\sqrt q}x^e\widehat b_e(x)\,dx,\\
 \partial_z\beta_e=-\frac{q_z}{2q}r^{e+1}b_e(r,z,t),\qquad
 \int_0^\infty r^e\partial_zb_e\,dr=0.
\end{gathered}
\end{equation}
The first identity is the product and chain rule. The second is the
same radial substitution as the normalization. Differentiating its
upper limit proves the third. Differentiating the normalization under
the uniformly compact radial integral proves the fourth. In particular
none of these identities assigns a sign to $q_z$.

\subsection{Corrected radial primitives and the complete symmetric tensor}

Define the actual corrected primitives
\begin{equation}\label{mcstress:sigma}
 \sigma_e(r,z,t)=-r^{-e}\int_0^r s^e
       \bigl(F_e(s,z,t)-\partial_z(b_e(s,z,t)J_e(z,t))\bigr)\,ds,
 \qquad e=1,2.
\end{equation}
Their exact radial equations are
\begin{equation}\label{mcstress:sigma_equation}
 (\partial_r+e/r)\sigma_e
      =-F_e+\partial_z(b_eJ_e).
\end{equation}
This follows by multiplying \eqref{mcstress:sigma} by $r^e$ and
using the fundamental theorem of calculus.

For comparison with the initially proposed primitive, put
\[
 \sigma_e^{\rm old}=-r^{-e}\int_0^r s^e
                       (F_e-b_e\partial_zJ_e)\,ds.
\]
Every missing term can be displayed exactly:
\begin{equation}\label{mcstress:primitive_correction}
 \sigma_e=\sigma_e^{\rm old}+r^{-e}J_e\partial_z\beta_e
         =\sigma_e^{\rm old}-\frac{rq_z}{2q}b_eJ_e.
\end{equation}
Indeed expanding $\partial_z(b_eJ_e)$ in
\eqref{mcstress:sigma} proves the first equality, and
\eqref{mcstress:bump_derivatives} proves the second. With the old
primitive the angular or axial stress divergence below would instead
be $-F_e-(\partial_zb_e)J_e$. Thus the correction is necessary for the
retained varying bump. It vanishes in the special case $q_z=0$.

There is an exact expression entirely in the original tensor entries:
\begin{equation}\label{mcstress:compact_flux_correction}
\begin{gathered}
 C_e(r,z,t)=\int_0^r s^eS_{z\tau_e}(s,z,t)\,ds
                                    -\beta_e(r,z,t)J_e(z,t),\\
 \sigma_e=-S_{r\tau_e}-r^{-e}\partial_zC_e.
\end{gathered}
\end{equation}
To prove it, integrate the corresponding component of
\eqref{mcstress:all_components} only up to $r$, obtaining
$\int_0^r s^eF_e\,ds=r^eS_{r\tau_e}
+\partial_z\int_0^r s^eS_{z\tau_e}\,ds$. Substitute this identity
and $\int_0^r s^e b_e\,ds=\beta_e$ into
\eqref{mcstress:sigma}. Below both radial supports $C_e=0$;
above both supports its two summands are $J_e$ and $J_e$, so it is
again zero. This proves its compact support as well as the exact
relation of the constructed radial stress to the actual covariance
tensor.

\begin{theorem}\label{mcstress:completion}
The real symmetric axisymmetric tensor with cylindrical entries
\begin{equation}\label{mcstress:tensor}
\begin{gathered}
 T_{rr}=0,\qquad T_{r\theta}=T_{\theta r}=\sigma_2,\qquad
 T_{rz}=T_{zr}=\sigma_1,\\
 T_{z\theta}=T_{\theta z}=-b_2J_2,\qquad
 T_{zz}=-b_1J_1,\qquad
 T_{\theta\theta}=rF_r+r\partial_z\sigma_1
\end{gathered}
\end{equation}
is smooth and compactly supported in an annular cylinder and obeys
the full Cartesian identity
\begin{equation}\label{mcstress:divergence}
 \operatorname{div}T=-(F_r e_r+F_2e_\theta+F_1e_z).
\end{equation}
If the term $rF_r$ is omitted only in $T_{\theta\theta}$, the exact
divergence is instead $-(0,F_2,F_1)$ in cylindrical components.
\end{theorem}
\begin{proof}
The integral of the numerator in \eqref{mcstress:sigma} over the
entire radial half-line is
\[
 M_e-\partial_z\left(J_e\int_0^\infty s^eb_e\,ds\right)
  =\partial_zJ_e-\partial_zJ_e=0.
\]
Thus the numerator primitive is zero below its support and zero above
it. Division by $r^e$ introduces no singularity because it is zero on
a fixed neighborhood of the axis. All derivatives under its fixed
compact radial integral exist, and the normalization identity holds
smoothly in $(z,t)$, so its extension by zero is smooth. Compact axial
support follows because both $F_e$ and $J_e$ vanish outside the fixed
compact axial set. The same statements hold for all entries of
\eqref{mcstress:tensor}. Passing to the Cartesian matrix
$O_\theta T O_\theta^T$ preserves smoothness away from the axis; near
the axis it is identically zero. Hence this is a smooth compactly
supported Cartesian tensor.

Apply the full symmetric axisymmetric divergence formula already
proved in \eqref{mcstress:all_components}. Its three rows are
\[
 \begin{aligned}
 (\operatorname{div}T)_r
   &=-\frac{rF_r+r\partial_z\sigma_1}{r}
                                     +\partial_z\sigma_1=-F_r,\\
 (\operatorname{div}T)_\theta
   &=(\partial_r+2/r)\sigma_2-\partial_z(b_2J_2)=-F_2,\\
 (\operatorname{div}T)_z
   &=(\partial_r+1/r)\sigma_1-\partial_z(b_1J_1)=-F_1.
 \end{aligned}
\]
The last two identities are \eqref{mcstress:sigma_equation}, including
every derivative of $b_e$. This proves the Cartesian vector identity
and, on omitting $rF_r$, its stated radial-zero variant.
\end{proof}

\subsection{Compact pressure, tracefree stress, and the exact gauge relation}

The trace and its prescribed pressure are explicitly
\begin{equation}\label{mcstress:pressure}
 \operatorname{tr}T=rF_r+r\partial_z\sigma_1-b_1J_1,
 \qquad
 p_T=\frac13\bigl(rF_r+r\partial_z\sigma_1-b_1J_1\bigr),
 \qquad T^0=T-p_TI_3.
\end{equation}
Then $\operatorname{tr}T^0=0$, and every field is smooth and compactly
supported in the same annular cylinder. The exact full cancellation is
\begin{equation}\label{mcstress:tracefree_cancellation}
 F+\operatorname{div}T^0+\nabla p_T=0.
\end{equation}
To check its sign, for any scalar $p$ Cartesian differentiation gives
$\operatorname{div}(pI_3)=\nabla p$. Therefore
$\operatorname{div}T^0=\operatorname{div}T-\nabla p_T
=-F-\nabla p_T$, which proves the displayed equation.
In components the new diagonal entries are
\[
 T^0_{rr}=-p_T,\quad
 T^0_{\theta\theta}=rF_r+r\partial_z\sigma_1-p_T,\quad
 T^0_{zz}=-b_1J_1-p_T.
\]
The axial stress divergence alone changes by $-\partial_zp_T$;
the pressure gradient contributes $+\partial_zp_T$. Thus the combined
stress--pressure correction preserves the exact tangential equations.
There is no claim that their stress term alone stays unchanged.

More generally $T-pI_3$ with any smooth compact scalar $p$ satisfies
$F+\operatorname{div}(T-pI_3)+\nabla p=0$.
The choice in \eqref{mcstress:pressure} is the unique one making
this stress tracefree. The map from symmetric tensors to
tracefree-tensor--scalar pairs,
\[
 T\longmapsto\bigl(T-\tfrac13\operatorname{tr}T\,I_3,
                         \tfrac13\operatorname{tr}T\bigr),
\]
has exact inverse $(U,p)\mapsto U+pI_3$ when $\operatorname{tr}U=0$.
It intertwines $\operatorname{div}$ with
$(U,p)\mapsto\operatorname{div}U+\nabla p$. This supplies the full
map between the two presentations, including compact support and signs.

\subsection{All total force and torque components and the divergence image}

The Cartesian components of the averaged force are
\[
 F=(F_r\cos\theta-F_2\sin\theta,
       F_r\sin\theta+F_2\cos\theta,F_1).
\]
Integration in $\theta$ with the physical volume element
$dx=r\,dr\,d\theta\,dz$ gives
\begin{equation}\label{mcstress:force_torque}
\begin{gathered}
 \int_{\mathbb R^3}F\,dx
  =2\pi e_z\int_{\mathbb R}M_1(z,t)\,dz=0,\\
 \int_{\mathbb R^3}x\times F\,dx
  =2\pi e_z\int_{\mathbb R}M_2(z,t)\,dz=0.
\end{gathered}
\end{equation}
For the torque calculation, direct multiplication in the oriented
basis gives
\[
 x\times F=-zF_2e_r+(zF_r-rF_1)e_\theta+rF_2e_z.
\]
The first two terms integrate to zero in $\theta$, and the third
gives the second equation of \eqref{mcstress:force_torque}.
Thus all three force and all three torque components vanish. The last
equalities use the already proved compact flux identities
\eqref{mcstress:flux_identity}. The same torque vanishes about every
fixed origin $x_0$, because subtracting $x_0$ changes it by
$-x_0\times\int F\,dx=0$. On a finite axial slab $a<z<b$, the
axial force and torque are respectively
$2\pi(J_1(b,t)-J_1(a,t))$ and
$2\pi(J_2(b,t)-J_2(a,t))$. The quantities need not vanish on an
individual slice.

There is also a direct Cartesian stress verification. Since
$F_j=-\partial_kT_{jk}$ and $T$ is compact, integration by parts gives
\begin{equation}\label{mcstress:first_moments}
 \int F_j\,dx=0,\qquad
 \int x_iF_j\,dx=\int T_{ji}\,dx=-\int S_{ji}\,dx.
\end{equation}
Symmetry makes the antisymmetric part of the first-moment matrix zero,
which is precisely all three torque components. Its symmetric part
is not asserted to vanish. The second equality with $S$ uses
$F=\operatorname{div}S$. Equivalently
$\operatorname{div}(T+S)=0$; multiplying this equation by $x_i$ and
integrating proves $\int(T+S)_{ji}\,dx=0$. In the tracefree-pressure
presentation the same first moment is
$\int T^0_{ji}\,dx+\delta_{ij}\int p_T\,dx$.
All these identities persist under any finite time differentiation on
a common compact support, by differentiation under the integral.

For completeness, this construction proves the exact image of compact
symmetric-stress divergence in the axisymmetric class. Let
an open axial interval $I_z$ contain the actual axial supports, with
$q>0$ smooth on a neighborhood of its relevant compact subsets at the
fixed time. In the following two spaces all supports are compact
subsets of $\{r>0,\ z\in I_z\}$. Thus the original bump is defined
everywhere it is used, and integrating between axial support endpoints
stays inside $I_z$.
Let
$\mathcal F_0$ be the space of smooth real axisymmetric compact forces
supported away from $r=0$ and satisfying exactly
\[
 \int_{\mathbb R}\int_0^\infty rF_1\,dr\,dz=0,
 \qquad
 \int_{\mathbb R}\int_0^\infty r^2F_2\,dr\,dz=0.
\]
Let $\mathcal S_c$ be the space of smooth real compact symmetric
axisymmetric Cartesian tensors supported away from $r=0$. Define
$\mathcal D:\mathcal S_c\to\mathcal F_0$ by
$\mathcal D(U)=-\operatorname{div}U$.
Integration by parts proves that its values satisfy both defining
identities of $\mathcal F_0$, and axisymmetry makes the other force
and torque components vanish as above.

For any $F\in\mathcal F_0$, compute its actual radial moments and set
\[
 J_e(z)=\int_{-\infty}^z\int_0^\infty r^eF_e(r,s)\,dr\,ds.
\]
This is a smooth compact axial function: its derivative is $M_e$, it
vanishes below the axial support, and the defining zero total integral
makes it vanish above that support. Use the fixed bumps
\eqref{mcstress:bumps}, the corrected primitive
\eqref{mcstress:sigma}, and all entries of
\eqref{mcstress:tensor}. The proof of Theorem~\ref{mcstress:completion}
uses only these identities and proves a compact tensor
$\mathcal C_bF$ with
\begin{equation}\label{mcstress:right_inverse}
 \mathcal D\mathcal C_bF=F\qquad(F\in\mathcal F_0).
\end{equation}
Thus the two total moments are necessary and sufficient, with an
explicit right inverse. For the actual force in
\eqref{mcstress:Ffull}, this primitive in $z$ equals the already
computed $J_e$ from $S$, because both have derivative $M_e$ and vanish
at the negative axial end. No additional compatibility condition has
been imposed on the actual candidate.

For fixed $q$ and bump profiles, all operations defining
$\mathcal C_b$ are linear. Hence
$\mathcal P_b=\mathcal C_b\mathcal D$ is a projection on
$\mathcal S_c$: \eqref{mcstress:right_inverse} gives
$\mathcal P_b^2=\mathcal P_b$ and
$\ker\mathcal P_b=\ker\mathcal D$.
In particular the completed actual tensor is
$T=\mathcal P_b(-S)$, and $T+S$ is a compact divergence-free symmetric
tensor. This gives an exact relation to the actual averaged covariance,
not merely equality of its two radial moments.

\subsection{Exact source chart, bump, stress, and pressure powers}

Fix one source chart and retain all of its constants:
\begin{equation}\label{mcstress:chart}
 r=\sqrt Q\,R,\qquad z=Q^D Z,\qquad1-t=QT,
 \qquad A=\tfrac12+h,\quad D=\tfrac12-h,\quad
 \varepsilon=Q^h.
\end{equation}
The dyadic label $Q$ is fixed under differentiation. The averaged quantities
are independent of $Y$ and $\theta$, so in this
subsection the complete normalized spatial derivatives are
$\partial_R$, $R^{-1}\partial_\theta$ and $\varepsilon\partial_Z$,
with the same cylindrical basis derivatives as before.
Define the exact chart representatives
\begin{equation}\label{mcstress:chart_representatives}
\begin{gathered}
 S^*=Q^{2A}S,\qquad F^*=Q^{2A+1/2}F,\qquad
 T^*=Q^{2A}T,\qquad \Sigma_e=Q^{2A}\sigma_e,\\
 J_e^*=Q^{2A-(e+1)/2}J_e
             =\int_0^\infty R^e S^*_{z\tau_e}\,dR,\qquad
 M_e^*=Q^{2A-e/2}M_e=\varepsilon\partial_ZJ_e^*,\\
 q_*=q/Q,\qquad
 b_e^*=Q^{(e+1)/2}b_e(\sqrt Q R,z,t)
       =q_*^{-(e+1)/2}\widehat b_e(R/\sqrt{q_*}).
\end{gathered}
\end{equation}
The two moment powers follow directly from
$r^e\,dr=Q^{(e+1)/2}R^e\,dR$ and
$\partial_z=Q^{-D}\partial_Z$; the remaining factor in $M_e^*$
is $Q^{1/2-D}=\varepsilon$. In particular the scaled bump satisfies
$\int R^e b_e^*\,dR=1$.

Every term of the corrected primitive has the same exact stress unit:
\begin{equation}\label{mcstress:chart_sigma}
\begin{gathered}
 \Sigma_e=-R^{-e}\int_0^R X^e
             \bigl(F_e^*-\varepsilon\partial_Z(b_e^*J_e^*)\bigr)\,dX,\\
 (\partial_R+e/R)\Sigma_e
       =-F_e^*+\varepsilon\partial_Z(b_e^*J_e^*),\\
 \Sigma_e=\Sigma_e^{\rm old}
       -\frac{\varepsilon R\partial_Zq_*}{2q_*}b_e^*J_e^*.
\end{gathered}
\end{equation}
To verify the first line, insert
$F_e=Q^{-2A-1/2}F_e^*$ and $b_eJ_e=Q^{-2A}b_e^*J_e^*$ in
\eqref{mcstress:sigma}, and change the integration variable to
$s=\sqrt Q X$. The second follows by differentiation. For the third,
the physical correction in \eqref{mcstress:primitive_correction}
has factor $\sqrt Q\,q_z/q$; using
$q_z=Q^{1-D}\partial_Zq_*$ gives
$\sqrt Q\,q_z/q=\varepsilon\partial_Zq_*/q_*$.
This verifies the original $q$ derivative, including its sign.
Also, with $C_e^*=Q^{2A-(e+1)/2}C_e$,
\[
 \Sigma_e=-S^*_{r\tau_e}
                     -\varepsilon R^{-e}\partial_Z C_e^*.
\]

The completed stress and its pressure become exactly
\begin{equation}\label{mcstress:chart_tensor}
\begin{gathered}
 T^*_{rr}=0,\quad T^*_{r\theta}=\Sigma_2,\quad
 T^*_{rz}=\Sigma_1,\quad T^*_{z\theta}=-b_2^*J_2^*,\quad
 T^*_{zz}=-b_1^*J_1^*,\\
 T^*_{\theta\theta}=RF_r^*+\varepsilon R\partial_Z\Sigma_1,
 \qquad
 p_T^*=Q^{2A}p_T
  =\tfrac13(RF_r^*+\varepsilon R\partial_Z\Sigma_1-b_1^*J_1^*),\\
 (T^0)^*=T^*-p_T^*I_3,\qquad
 F^*+\operatorname{div}_*(T^0)^*+\nabla_*p_T^*=0.
\end{gathered}
\end{equation}
For example the radial row is
$-T^*_{\theta\theta}/R+\varepsilon\partial_Z\Sigma_1=-F_r^*$,
and the tangential rows use \eqref{mcstress:chart_sigma}. The pressure
identity follows from the same Cartesian isotropic-tensor calculation
as before. Physical divergence and pressure gradient both have factor
$Q^{-2A-1/2}$ because $\nabla_x=Q^{-1/2}\nabla_*$ on the chart.
Thus the radial residual, hoop stress, both surviving axial fluxes,
varying-bump correction and compact pressure all have their exact
source units. No decay or uniform estimate at a singular endpoint is
inferred from these finite compact-support identities.


\section{Original-scale derivative costs of the actual compatibility defects}
\label{mcq:section}

The state and increment in this section are exactly the finite source state
and actual third-growth curl constructed in the preceding reader.
Write \(u_*=\langle u_*\rangle_\theta+w_*\), \(v=w_y\), and
\[
 C=\langle w_*\otimes v+v\otimes w_*+v\otimes v\rangle,
\]
where the outer brackets are the normalized angular and auxiliary
averages before physical evaluation. Every component below is in the
original oriented cylindrical orthonormal basis. The full moment
completion and source pressure calculation prove the formulas used here;
we also derive the scalar compatibility formulas directly to fix all signs.
Throughout this section the unadorned symbols \(P,J_\theta,J_z\)
denote the defect increments \(\Delta P,\Delta J_\theta,\Delta J_z\)
caused by this actual wave addition. They are not the full old-plus-new
source defects. The unchanged old defects must be added when applying
the source correction to the whole updated state.

\subsection{Actual mixed derivatives of the third-growth curl}

Retain the fixed finite set of source labels \(j\), the actual clocks
\(t_{c,j}=t_b+\alpha_jv_j\), and the two full shape fields \(W_{ja}\).
The original data \(t_b,r_a,S_3,u_i\) are independent of these spatial
labels. In particular
\[
 v=\sum_{j,a}Q_j^{-A}y_a(t_{c,j})W_{ja},\quad
 A=\tfrac12+h,\quad D=\tfrac12-h,\quad
 \partial_t t_{c,j}=\alpha_jQ_j^{-1-h},\quad \partial_z t_{c,j}=0 .
\]
Let \(Y_{1,n}=T_n\) and \(Y_{2,n}=O_n\) be the complete finite
derivative arrays already proved for the actual coupled solution.
Thus \(Y_{1,0}=(A_0/\mu)\widehat\lambda_3^{-7/8}\),
\(Y_{2,0}=2A_0\widehat\lambda_3^{1/8}/(\sigma_2\sqrt{n_i})\),
\(Y_{1,1}=2\gamma_iY_{1,0}\), and
\(Y_{2,1}=\gamma_iA_0\widehat\lambda_3^{1/8}/(\sigma_2\sqrt{n_i})\).
Higher entries are those closed recurrences in the retained moving-parent
data; no independent amplitude bound is introduced here.

For a pair \(I=(a,b)\in\mathbb N^2\), put
\(\partial^I=\partial_z^a\partial_t^b\) and
\(I_z=I+(1,0)\). On a fixed chart the full time operator is
\(\mathfrak t_{*,j}=-\varepsilon_j\partial_{T_j}+c_{i,j}N_{i,j}\).
The exact mixed derivative formula is
\begin{equation}\label{mcq:actual_derivative}
\begin{split}
 \partial_z^a\partial_t^b v
 =\sum_{j,c}\sum_{s=0}^b&
 \binom bs Q_j^{-A-aD-b(1+h)}\alpha_j^s
 y_c^{(s)}(t_{c,j})\\
 &\hspace{4mm}\cdot
 \partial_{Z_j}^a\mathfrak t_{*,j}^{\,b-s}W_{jc}.
\end{split}
\end{equation}
All operators on the shapes include their phase, frame and cutoff
derivatives. The shape time operator is the full derivative after
auxiliary evaluation, not a derivative that holds the fast phase fixed.
To prove the identity, differentiate the physical product \(b\) times
in time. If \(s\) of the derivatives act on \(y_c\), its factor is
\((\alpha_jQ_j^{-1-h})^sy_c^{(s)}\); the other \(b-s\) derivatives
give \(Q_j^{-(b-s)(1+h)}\mathfrak t_{*,j}^{b-s}W_{jc}\).
There are \(\binom bs\) such choices. The \(a\) axial derivatives act
only on the shape and give \(Q_j^{-aD}\partial_{Z_j}^a\).
The axial derivative commutes with the stated time operator, since
its coefficients are fixed label constants. Combining the powers proves
\eqref{mcq:actual_derivative} without dropping any derivative.

For each \(e\in\{-1,0,1,2\}\), define the precise weighted norm
\[
 \|f\|_{[e]}^2=\int_0^\infty r^e\langle |f|^2\rangle\,dr,\qquad
 \|F\|_{[e],j}^2=\int_0^\infty R_j^e\langle |F|^2\rangle\,dR_j.
\]
Every field being integrated is supported in a compact annulus away
from zero, so the negative weight introduces no singular integral.
Define the explicit actual-shape budget
\begin{equation}\label{mcq:V}
\begin{split}
 \mathcal V_{e,(a,b)}
 =\sum_{j,c}\sum_{s=0}^b&
 \binom bs \alpha_j^sY_{c,s}
 Q_j^{(e+1)/4-A-aD-b(1+h)}\\
 &\hspace{4mm}\cdot
 \|\partial_{Z_j}^a\mathfrak t_{*,j}^{\,b-s}W_{jc}\|_{[e],j}.
\end{split}
\end{equation}
Then \(\|\partial^I v\|_{[e]}\le\mathcal V_{e,I}\).
Indeed the measure identity
\(r^e\,dr=Q_j^{(e+1)/2}R_j^e\,dR_j\), valid also for \(e=-1\),
and the triangle inequality applied to \eqref{mcq:actual_derivative}
give this bound. There is no inter-label orthogonality assumption.
Put \(\mathcal W_{e,I}=\|\partial^I w_*\|_{[e]}\), restricted to a
fixed radial set enclosing the supports of the actual new fields
and all derivatives at issue. The old wave includes its existing curl
corrections. These are norms of the specified source state, not missing
amplitude variables.

\subsection{A contraction estimate that retains every product derivative}

For \(J=(j_1,j_2)\le I=(a,b)\), set
\(\binom IJ=\binom a{j_1}\binom b{j_2}\), and define
\begin{equation}\label{mcq:B}
\begin{split}
 \mathcal B_{e,I}=\sum_{J\le I}\binom IJ
 \bigl(&\mathcal W_{e,J}\mathcal V_{e,I-J}
       +\mathcal V_{e,J}\mathcal W_{e,I-J}\\
       &+\mathcal V_{e,J}\mathcal V_{e,I-J}\bigr).
\end{split}
\end{equation}
If \(E\) is any fixed real symmetric three-by-three matrix in
cylindrical components, then
\begin{equation}\label{mcq:contraction}
 \left|\partial^I\int r^e(E:C)\,dr\right|
       \le \|E\|_{\mathrm{op}}\mathcal B_{e,I}.
\end{equation}
For a single outer product,
\(E:(f\otimes g)=f^TEg\), whose absolute value is at most
\(\|E\|_{\mathrm{op}}|f||g|\). The repeated product rule gives exactly
the sum over \(J\le I\) in \eqref{mcq:B}. Cauchy--Schwarz on the
original weighted angular-auxiliary-radial measure gives each displayed
product of norms. Compact smooth support justifies differentiation
under the integral. Axial and time derivatives do not differentiate
the cylindrical basis. This proves \eqref{mcq:contraction} for every
specified order.
For the time derivative, the full lifted physical operator is
\(\mathcal D_t=\partial_t+N_{\rm abs}\), with \(N_{\rm abs}\) a
constant torus direction. The Haar integral of \(N_{\rm abs}f\) is
zero by periodicity. Therefore
\(\partial_t\langle f\rangle=\langle\mathcal D_tf\rangle\).
The same identity applies repeatedly, and axial differentiation commutes
with this average and this constant torus direction. Consequently the
ordinary mixed derivatives of the averaged tensor in
\eqref{mcq:contraction} are exactly the averages of the full physical
derivatives estimated in \eqref{mcq:actual_derivative}.

\subsection{The actual three defects and the pressure-corrected axial flux}

Since the wave increment has zero angular mean and is divergence-free,
the change of the radial mean pressure source is
\[
 g=-(\partial_r+r^{-1})C_{rr}-\partial_zC_{zr}+r^{-1}C_{\theta\theta}.
\]
The three actual source compatibility-defect increments are therefore
\begin{equation}\label{mcq:defects}
\begin{split}
 P&=\int\frac{C_{\theta\theta}-C_{rr}}r\,dr
                    -\partial_z\int C_{zr}\,dr,\\
 J_\theta&=\int r^2C_{z\theta}\,dr,\\
 J_z&=\int r\left(C_{zz}-\tfrac12C_{rr}
                              -\tfrac12C_{\theta\theta}\right)\,dr
                  +\tfrac12\partial_z\int r^2C_{zr}\,dr .
\end{split}
\end{equation}
The first formula is \(\int g\,dr\), using the zero boundary integral
of \(\partial_r C_{rr}\). For the last, integration by parts gives
\[
 \int r^2 g\,dr
  =\int r(C_{rr}+C_{\theta\theta})\,dr
                      -\partial_z\int r^2C_{zr}\,dr .
\]
Substitute it into the source definition
\(J_z=\int rC_{zz}\,dr-\frac12\int r^2g\,dr\).
This proves all three identities in \eqref{mcq:defects}.
In particular the axial source defect is not replaced by the raw flux
\(\int rC_{zz}\,dr\).

Let the source pressure bump be the original
\[
 \rho(r,z,t)=q^{-1/2}\widehat\rho(r/\sqrt q),\qquad
 \int\widehat\rho(x)\,dx=1,\qquad
 c_0=\tfrac12\int x^2\widehat\rho(x)\,dx .
\]
It has \(\int\rho\,dr=1\) and
\(c_\rho=\frac12\int r^2\rho\,dr=c_0q\).
The auxiliary average of the compact source pressure increment is
\[
 \overline{\Delta p_m}(r)=\int_0^r(g(s)-\rho(s)P)\,ds .
\]
Its derivative is \(g-\rho P\), and its integrand has zero total
radial integral, giving compact support. Its weighted pressure moment is
\(\int r\overline{\Delta p_m}\,dr=-\frac12\int r^2g\,dr+c_\rho P\).
Here is the exact map to the full auxiliary-dependent source operation.
Let \(\Delta g_r\) be the angularly averaged source increment before
auxiliary averaging, so that \(\langle\Delta g_r\rangle_Y=g\), and
put \(f=\Delta g_r-\rho P\). The source's shifted primitive uses
\[
 \begin{split}
 \mathcal I f(r,Y)&=\int_0^r
       f(s,Y+(s^{d_r}-r^{d_r})v_r)\,ds,\\
 \mathcal J f(r,Y)&=\int_0^\infty
       f(s,Y+(s^{d_r}-r^{d_r})v_r)\,ds,\qquad
 \mathcal I_cf=\mathcal If-\chi_m\mathcal Jf .
 \end{split}
\]
The fixed source exponent and direction are the original \(d_r,v_r\),
and \(\chi_m\) is its inner-zero, outer-one radial cutoff. The full
pressure increment is \(\Delta p_m=\mathcal I_cf\).
Haar translation invariance gives
\(\langle\mathcal Jf\rangle_Y=\int(g-\rho P)\,dr=0\) and
\(\langle\mathcal If\rangle_Y=\int_0^r(g-\rho P)\,ds\).
This proves the asserted averaged pressure formula from the full map.
The lifted radial derivative cancels the shift derivative and yields
\[
 \mathcal D_r\Delta p_m-\Delta g_r
       =-\rho P-(\partial_r\chi_m)\mathcal Jf .
\]
The second term has zero auxiliary average but need not vanish before
physical evaluation. No estimate for its singular-time flatness is
inferred from its zero average.
Consequently the double-averaged axial residual-moment increment after
this reconstruction is
\[
 \partial_zK_z,\qquad K_z=J_z+c_0qP .
\]
The full product, including the derivatives of \(q\), is retained.

The following simultaneous quantitative estimates hold at every finite
mixed order:
\begin{equation}\label{mcq:defect_bounds}
\begin{split}
 |\partial^IP|&\le
       \mathcal B_{-1,I}+\tfrac12\mathcal B_{0,I_z},\\
 |\partial^IJ_\theta|&\le\tfrac12\mathcal B_{2,I},\\
 |\partial^IJ_z|&\le
       \mathcal B_{1,I}+\tfrac14\mathcal B_{2,I_z},\\
 |\partial^IK_z|&\le
       \mathcal B_{1,I}+\tfrac14\mathcal B_{2,I_z}\\
 &\quad+|c_0|\sum_{J\le I}\binom IJ
       |\partial^Jq|\left(
       \mathcal B_{-1,I-J}+\tfrac12\mathcal B_{0,(I-J)_z}\right).
\end{split}
\end{equation}
For the first line the diagonal contraction matrix is
\(\operatorname{diag}(-1,1,0)\), of norm one, whereas the \(rz\)
entry is the contraction against the symmetric matrix with \(rz,zr\)
entries \(1/2\), of norm \(1/2\). For the second line use the
analogous \(z\theta\) matrix. For the third the diagonal matrix is
\(\operatorname{diag}(-1/2,-1/2,1)\), of norm one, and the
coefficient \(1/2\) multiplying \(C_{zr}\) contributes an additional
factor \(1/2\). These facts and \eqref{mcq:contraction} prove the first
three bounds. The exact product rule for \(qP\) proves the fourth.
Replacing \(I\) by \(I_z\) gives bounds for the two residual moments
that actually enter the source radial inverse. Thus every pressure,
radial and axial coupling has a bound in the original shape fields
and third-growth constants.

\subsection{Exact derivatives of the original moving scale and bumps}

The source scale is determined by
\[
 q-z^2q^{2h}=1-t,\qquad z=q^D\eta,\qquad
 D=\tfrac12-h,\qquad 0<h<1/100,\quad -1<\eta<1 .
\]
Set \(B(\eta)=1-2h\eta^2>1-2h>0\). Implicit differentiation of the
first equation, followed by the displayed formula for \(z\), gives
\begin{equation}\label{mcq:q_derivatives}
\begin{gathered}
 q_z=\frac{2\eta q^A}{B},\qquad q_t=-\frac1B,\\
 \eta_z=\frac{q^{-D}(1-\eta^2)}B,\qquad
 \eta_t=\frac{D\eta q^{-1}}B .
\end{gathered}
\end{equation}
For example \(q_z=2zq^{2h}/B\), and \(D+2h=A\).
Differentiating \(\eta=zq^{-D}\) gives the last two formulas;
the identity \(2D+2h=1\) gives the factor \(1-\eta^2\).
This proves the derivative dictionary without assigning a sign to \(q_z\).

Put \(x=r/\sqrt q\). For a smooth function \(f(x,\eta)\) define the
explicit differential expressions
\[
 \mathcal Z_\sigma f=
 \frac{2\sigma\eta f-\eta x f_x+(1-\eta^2)f_\eta}{B},\qquad
 \mathcal T_\sigma f=
 \frac{-\sigma f+\frac12x f_x+D\eta f_\eta}{B}.
\]
Equations \eqref{mcq:q_derivatives} and the ordinary chain rule prove
\begin{equation}\label{mcq:dilation_dictionary}
 \partial_z(q^\sigma f)=q^{\sigma-D}\mathcal Z_\sigma f,\quad
 \partial_t(q^\sigma f)=q^{\sigma-1}\mathcal T_\sigma f,\quad
 \partial_r(q^\sigma f)=q^{\sigma-1/2}f_x .
\end{equation}
For the last formula \(q,\eta\) are independent of \(r\).
For the first two, differentiation of \(x\) gives respectively
\(-\eta xq^{-D}/B\) and \(xq^{-1}/(2B)\), proving every term.

In particular let the original radial moment bump be
\[
 b_e=q^{-(e+1)/2}\widehat b_e(x),\qquad
 \int x^e\widehat b_e(x)\,dx=1 .
\]
For any \(k,a,b\ge0\), its derivative is exactly
\begin{equation}\label{mcq:bump_derivatives}
 \partial_r^k\partial_z^a\partial_t^b b_e
 =q^{-(e+1+k)/2-aD-b} f_{e;k,a,b}(x,\eta).
\end{equation}
Here the coefficient is given by a finite recursion: start at
\(\sigma=-(e+1+k)/2\) and \(f=\widehat b_e^{(k)}\);
repeat \(f\leftarrow\mathcal Z_\sigma f\),
\(\sigma\leftarrow\sigma-D\), \(a\) times, and then
\(f\leftarrow\mathcal T_\sigma f\),
\(\sigma\leftarrow\sigma-1\), \(b\) times.
The result is \(f_{e;k,a,b}\). This proves the identity by induction
using \eqref{mcq:dilation_dictionary}. Smooth mixed derivatives commute,
so this fixed order computes the stated derivative.
Every coefficient is explicitly a finite rational expression in \(x,\eta\),
the retained constants, and derivatives of the specified bump, with
nonzero denominator a power of \(B\). On its compact \(x\) support
and \(|\eta|\le1\) its maximum is finite and computable directly from
that expression. Thus \eqref{mcq:bump_derivatives} supplies the original
power of \(q\) at every order.
The pressure bump \(\rho\) uses exactly this recursion with \(e=0\).
Likewise every derivative in \eqref{mcq:defect_bounds} is explicit:
\(\partial_z^a\partial_t^bq=q^{1-aD-b}Q_{a,b}(\eta)\), where the same
recursion starts with \(\sigma=1,f=1\).

For completeness, the weighted cumulative bump has the exact form
\[
 \beta_e(r,z,t)=\int_0^rs^eb_e(s,z,t)\,ds
              =\int_0^{r/\sqrt q}x^e\widehat b_e(x)\,dx .
\]
Its axial derivative is
\begin{equation}\label{mcq:bump_movement}
 \partial_z\beta_e
 =-\frac{q_z}{2q}\,r^{e+1}b_e .
\end{equation}
This follows by differentiating its upper endpoint, which is
\(-rq_z/(2q^{3/2})\). Inserting the definition of \(b_e\) gives
the displayed original physical factors. Therefore replacing the
primitive of \(F_e-b_e\partial_zJ_e\) by the exact primitive of
\(F_e-\partial_z(b_eJ_e)\) changes the stress coefficient by
\[
 -\frac{rq_z}{2q}\,b_eJ_e .
\]
This is the moving-bump term required in the complete stress construction.
Its sign and all further derivatives are fixed by
\eqref{mcq:q_derivatives}--\eqref{mcq:bump_movement}; it is not set to zero
by holding the physical source scale fixed during axial differentiation.

These calculations supply finite quantitative inputs for the actual
source pressure and moment operations. They do not identify the
distinct direct candidate with a completed infinite source correction
sequence: the complete source-operation formulas retain the new
velocity, covariance, pressure and nonlinear defects that such a step
actually creates.


\section{Physical derivative and full vector-force costs of the five-row mean operation}
\label{mcmean:section}

This section completes the finite mean operation on the actual current
source state. All coordinates are the original right-handed cylindrical
coordinates. The physical viscosity is the fixed number
\(\nu_{\rm NS}>0\). Write the full current velocity as
\(U=(b+\beta,V+v,G+\gamma)+w\), where \(w\) includes every current
wave and curl correction. During the operation below this wave is held
fixed. The three targets are the full current defects
\begin{equation}\label{mcmean:full_targets}
 \mathcal P=P_{\rm old}+\Delta P,\qquad
 \mathcal J_\theta=(J_\theta)_{\rm old}+\Delta J_\theta,\qquad
 \mathcal J_z=(J_z)_{\rm old}+\Delta J_z.
\end{equation}
Here the increments are the pressure-corrected actual wave increments
of \eqref{mcq:defects}. If other finite source operations have already
been performed, ``old'' means their actual entering state and the
increments are recomputed for the actual addition. Thus none of the
following formulas substitutes an increment for a full target.

Retain, without a change of scale or profile constant,
\begin{equation}\label{mcmean:original_scale}
\begin{gathered}
 q-z^2q^{2h}=1-t,\quad z=q^D\eta,\quad x=r/\sqrt q,\quad
 A=\tfrac12+h,\quad D=\tfrac12-h,\\
 0<h<1/100,\quad q>0,\quad -1<\eta<1,\quad
 B(\eta)=1-2h\eta^2,\\
 V=q^{-A}a(\eta)x^{-1-2\lambda},\quad G=0
 \quad\hbox{on the reserved mean patch},\\
 a(\eta)=\frac{2^{1/2+\lambda}c_{\rm patch}}{1+\eta^2},
 \qquad \lambda>0,\qquad |a(\eta)|\ge a_0>0.
\end{gathered}
\end{equation}
The symbol \(a(\eta)\) in this display is a scalar base coefficient;
below the bold symbol \(\boldsymbol a\) denotes the vector increment.
The original constant \(c_{\rm patch}\) can have either nonzero sign.

\subsection{Exact inverse profiles and the five physical rows}

Use precisely the bumps and spacing of the five-row operation:
\[
 \varphi_j(x)=e^{-jd_b}\varphi_0(e^{-jd_b}x),\qquad
 \mu_p=\int_0^\infty x^p\varphi_0(x)\,dx>0,
 \qquad t_p=e^{d_bp},\quad d_b>0.
\]
The nonnegative nonzero smooth bump \(\varphi_0\) and its three
dilates have support in the reserved open interval \(I_m\Subset(0,\infty)\).
The substitution \(x=e^{jd_b}\xi\) proves exactly
\(\int x^p\varphi_j(x)\,dx=\mu_p t_p^j\).
Number matrix rows and columns starting at zero, and define
\begin{equation}\label{mcmean:vandermonde}
\begin{gathered}
 V_\theta=[t_p^j]_{p=2,-2-2\lambda,-2\lambda; j=0,1,2},
 \qquad B_\theta=V_\theta^{-1},\\
 V_z=[t_p^j]_{p=1,1-2\lambda; j=0,1},
 \qquad B_z=V_z^{-1},\\
 A_\theta=\operatorname{diag}
  (\mu_2,2a\mu_{-2-2\lambda},-a\mu_{-2\lambda})V_\theta,
 \qquad
 A_z=\operatorname{diag}(\mu_1,a\mu_{1-2\lambda})V_z.
\end{gathered}
\end{equation}
For three row values \(t_0,t_1,t_2\), subtraction of the first row
from each later row, followed by factoring, gives
\(\det V_\theta=(t_1-t_0)(t_2-t_0)(t_2-t_1)\).
For two row values the determinant is \(t_1-t_0\).
The original exponents in each displayed matrix are distinct because
\(\lambda>0\), and the exponential with \(d_b>0\) is injective.
The row factors are nonzero. This proves both inverses exist at the
actual fixed source parameters. Explicitly each entry of either
\(B\) is its corresponding transposed cofactor divided by the stated
determinant; this is a finite formula in the original exponentials.

Define the following actual profiles, including their \(\eta\) dependence:
\begin{equation}\label{mcmean:profiles}
\begin{split}
 V_P(x,\eta)&=-\frac{1}{2a(\eta)\mu_{-2-2\lambda}}
                   \sum_{j=0}^2B_\theta[j,1]\varphi_j(x),\\
 V_{J_z}(x,\eta)&=\frac{1}{a(\eta)\mu_{-2\lambda}}
                   \sum_{j=0}^2B_\theta[j,2]\varphi_j(x),\\
 G_{J_\theta}(x,\eta)&=-\frac{1}{a(\eta)\mu_{1-2\lambda}}
                   \sum_{j=0}^1B_z[j,1]\varphi_j(x).
\end{split}
\end{equation}
The signs, the factor two, and all three negative powers in the moment
denominators are part of the formulas. Multiplying the columns by
the exact rows in \eqref{mcmean:vandermonde} proves the eight scalar
profile identities
\begin{equation}\label{mcmean:profile_moments}
\begin{array}{c|cc}
 &V_P&V_{J_z}\\ \hline
 \int x^2(\cdot)\,dx&0&0\\
 2a\int x^{-2-2\lambda}(\cdot)\,dx&-1&0\\
 -a\int x^{-2\lambda}(\cdot)\,dx&0&-1
\end{array}
\qquad
\begin{array}{c|c}
 &G_{J_\theta}\\ \hline
 \int x(\cdot)\,dx&0\\
 a\int x^{1-2\lambda}(\cdot)\,dx&-1
\end{array}.
\end{equation}
There are eight entries in these two tables; each follows from
\(V_\theta B_\theta=I_3\) or \(V_zB_z=I_2\), with the row factor
left in place. The complete matrix identities include all nine and
all four scalar inverse identities, respectively.

The source coefficients are
\[
 \boldsymbol u=A_\theta^{-1}
       (0,-q^{2A}\mathcal P,-q^{2A-1}\mathcal J_z)^T,
 \qquad
 \boldsymbol s=A_z^{-1}
       (0,-q^{2A-3/2}\mathcal J_\theta)^T.
\]
Since \((DV)^{-1}=V^{-1}D^{-1}\), substitution of these coefficients
in the original physical increments gives
\begin{equation}\label{mcmean:tangential_increment}
 \delta v=q^A\mathcal P V_P+q^{A-1}\mathcal J_zV_{J_z},
 \qquad
 \delta\gamma=q^{A-3/2}\mathcal J_\theta G_{J_\theta}.
\end{equation}
For example the fifth-row coefficient in \(\boldsymbol u\) is
\((-q^{2A-1}\mathcal J_z)/(-a\mu_{-2\lambda})\), which proves
the positive sign in \(V_{J_z}\). Multiplication by the original
outer factor \(q^{-A}\) proves its exponent \(A-1\).
The other two terms follow by the same displayed diagonal products.

Changing variables \(r=\sqrt q\,x\) in each integral, and using
\eqref{mcmean:profile_moments}, now proves the full five equations
\begin{equation}\label{mcmean:five_rows}
\begin{gathered}
 \int r^2\delta v\,dr=0,\qquad
 \int r\delta\gamma\,dr=0,\qquad
 \int\frac{2V}{r}\delta v\,dr=-\mathcal P,\\
 \int r^2(G\delta v+V\delta\gamma)\,dr=-\mathcal J_\theta,\\
 \int(2rG\delta\gamma-rV\delta v)\,dr=-\mathcal J_z.
\end{gathered}
\end{equation}
Indeed the pressure term with \(V_P\) has scale
\(q^{-A}q^A=1\); the angular flux with \(G_{J_\theta}\) has scale
\(q^{3/2}q^{-A}q^{A-3/2}=1\); and the fifth row with
\(V_{J_z}\) has scale \(q q^{-A}q^{A-1}=1\).
Every other term is one of the zero entries in the tables. The two
homogeneous velocity moments use the first rows of those tables.
This verifies the physical equations, including the single factor
\(r\) multiplying \(V\delta v\) in the fifth row.

\subsection{Compact potential, induced radial velocity, and support}

Define the combined potential profile
\begin{equation}\label{mcmean:H}
 H_\theta(x,\eta)=\frac1x\int_0^x\xi G_{J_\theta}(\xi,\eta)\,d\xi.
\end{equation}
Choose \(0<x_-<x_+\) such that the convex hull of all the bump
supports lies in \([x_-,x_+]\Subset I_m\). For \(x<x_-\) the
integral in \eqref{mcmean:H} is zero. For \(x>x_+\) it is the
zero axial moment in \eqref{mcmean:profile_moments}.
It follows that \(H_\theta\), together with every derivative, is
supported in this fixed compact interval. Its smoothness at the two
collars follows from the fact that the integrand is identically zero
there and that its total integral is zero. In particular it extends
smoothly by zero to \(x\le0\).

The primitive of an individual \(\varphi_j\) is
\(\psi_j=x^{-1}\int_0^x\xi\varphi_j(\xi)\,d\xi\), which equals
\(\mu_1e^{jd_b}/x\) past its support. Thus an individual primitive
can have a tail. The exact identity
\(\sum_{j=0}^1 B_z[j,1]e^{jd_b}=0\) cancels these tails in
\(H_\theta=-(a\mu_{1-2\lambda})^{-1}\sum_j B_z[j,1]\psi_j\).
The identity holds for every \(\eta\), so its \(\eta\) derivatives
also vanish. This proves compactness for the combined profile and
all derivatives without assuming compactness for the summands.

The physical potential and radial increment are exactly
\begin{equation}\label{mcmean:potential_radial}
\begin{split}
 \Psi(r,z,t)&=\frac1r\int_0^r s\delta\gamma(s,z,t)\,ds
             =q^{A-1}\mathcal J_\theta H_\theta(x,\eta),\\
 \delta\beta&=-\partial_z\Psi\\
 &=-q^{A-1}(\partial_z\mathcal J_\theta)H_\theta
   -q^{A-1-D}\mathcal J_\theta\mathcal Z_{A-1}H_\theta.
\end{split}
\end{equation}
The substitution in the integral supplies one factor \(q\) and its
outer \(1/r\) supplies \(q^{-1/2}\), giving the exponent \(A-1\).
The operator in the second line is explicitly
\[
 \mathcal Z_\sigma f
 =\frac{2\sigma\eta f-\eta x f_x+(1-\eta^2)f_\eta}{B(\eta)}.
\]
The chain rule \eqref{mcq:dilation_dictionary} proves the last line,
including the \(\eta\) derivative of the original factor \(1/a\).
Direct differentiation of \(xH_\theta\) gives
\((\partial_x+x^{-1})H_\theta=G_{J_\theta}\), so
\((\partial_r+r^{-1})\Psi=\delta\gamma\). Commuting the smooth
physical \(r,z\) derivatives proves
\begin{equation}\label{mcmean:divergence}
 (\partial_r+r^{-1})\delta\beta+\partial_z\delta\gamma=0.
\end{equation}
The three velocity increments and \(\Psi\) are auxiliary independent
and angularly independent; their support lies in
\(\sqrt q[x_-,x_+]\), including any intervals between the bumps.

\subsection{Every mixed physical derivative, with original scale powers}

For clarity the original differential dictionary is
\[
 q_z=\frac{2\eta q^A}{B},\quad q_t=-\frac1B,\quad
 \eta_z=\frac{q^{-D}(1-\eta^2)}B,\quad
 \eta_t=\frac{D\eta q^{-1}}B,
 \qquad
 \mathcal T_\sigma f=
 \frac{-\sigma f+\tfrac12x f_x+D\eta f_\eta}{B}.
\]
These equations follow by implicit differentiation of
\eqref{mcmean:original_scale}; in particular
\(1-2hz^2q^{2h-1}=B\) and \(D+2h=A\).
Differentiating \(x=rq^{-1/2}\) supplies
\(x_z=-\eta xq^{-D}/B\) and \(x_t=xq^{-1}/(2B)\).
Consequently
\[
 \partial_r(q^\sigma F)=q^{\sigma-1/2}F_x,\quad
 \partial_z(q^\sigma F)=q^{\sigma-D}\mathcal Z_\sigma F,\quad
 \partial_t(q^\sigma F)=q^{\sigma-1}\mathcal T_\sigma F.
\]
Here \(F=F(x,\eta)\). These are physical derivatives of the moving
scale; the time loss is the original one in this scale.

For \(k,m,n\ge0\) define \(F_{\sigma;k,m,n}\) by the following
finite recursion. Start with \(f=\partial_x^kF\) and
\(s=\sigma-k/2\). Repeat
\(f\leftarrow\mathcal Z_s f,\ s\leftarrow s-D\) exactly \(m\)
times; then repeat \(f\leftarrow\mathcal T_s f,\ s\leftarrow s-1\)
exactly \(n\) times. The final \(f\) is \(F_{\sigma;k,m,n}\).
For any actual scalar target \(X=X(z,t)\), the exact product formula is
\begin{equation}\label{mcmean:all_derivatives}
\begin{split}
 \partial_r^k\partial_z^m\partial_t^n(q^\sigma XF)
  =\sum_{i=0}^m\sum_{j=0}^n&\binom mi\binom nj
     (\partial_z^i\partial_t^jX)\\
  &\cdot q^{\sigma-k/2-(m-i)D-(n-j)}
       F_{\sigma;k,m-i,n-j}(x,\eta).
\end{split}
\end{equation}
To prove it, first apply \(k\) radial derivatives, which do not act
on \(X,q,\eta\). The repeated product rule for the \(m\) axial and
\(n\) time derivatives assigns \(i,j\) derivatives to \(X\), with
the two binomial multiplicities displayed. Apply the three chain
rules above successively to the remaining factor. Each axial
derivative lowers its current exponent by \(D\), and each time
derivative lowers it by one, while its current exponent is the index
in the next operator. This is precisely the stated recursion.
Smooth mixed derivatives commute, so the chosen radial-then-axial-
then-time order computes the derivative written in the formula.

This proves every mixed derivative of \(\delta v,\delta\gamma,\Psi\)
by using, respectively, the two terms and the one term in
\eqref{mcmean:tangential_increment} and \eqref{mcmean:potential_radial}.
It proves every derivative of \(\delta\beta\) either by applying
\eqref{mcmean:all_derivatives} to its two displayed terms or by
\begin{equation}\label{mcmean:radial_all_derivatives}
 \partial_r^k\partial_z^m\partial_t^n\delta\beta
 =-\partial_r^k\partial_z^{m+1}\partial_t^n
          (q^{A-1}\mathcal J_\theta H_\theta).
\end{equation}
This equality retains, in particular, the derivative of the full
target and the additional original factor \(q^{-D}\).

The recursion is an explicit finite rational expression in
\(x,\eta,h\), the retained profile derivatives, and powers of
\(B^{-1}\). The other profile constants are exactly the inverse
Vandermonde entries and moments in \eqref{mcmean:profiles}.
In fact the reciprocal base coefficient has the exact derivatives
\begin{equation}\label{mcmean:reciprocal_a}
 \partial_\eta^s(a^{-1})=
 \frac{1}{2^{1/2+\lambda}c_{\rm patch}}
 \begin{cases}1+\eta^2&s=0,\\2\eta&s=1,\\2&s=2,\\0&s\ge3.
 \end{cases}
\end{equation}
Also \(\varphi_j^{(k)}(x)=e^{-(k+1)jd_b}
\varphi_0^{(k)}(e^{-jd_b}x)\).
These two formulas give every mixed derivative of the three profiles
without an unspecified inverse-matrix derivative constant.

For a similarly explicit derivative of \(H_\theta\), write
\(F_0(x,\eta)=\int_0^x\xi G_{J_\theta}(\xi,\eta)\,d\xi\).
For \(\ell\ge1\) the fundamental theorem of calculus and the
product rule give
\[
 \partial_x^\ell F_0=x\partial_x^{\ell-1}G_{J_\theta}
              +(\ell-1)\partial_x^{\ell-2}G_{J_\theta},
\]
where the second term is defined to be zero for \(\ell=1\).
Therefore, for \(x>0\),
\begin{equation}\label{mcmean:H_derivatives}
 \partial_x^k\partial_\eta^sH_\theta
 =\sum_{\ell=0}^k\binom k\ell(-1)^{k-\ell}(k-\ell)!
    x^{-1-k+\ell}\partial_x^\ell\partial_\eta^sF_0.
\end{equation}
The \(\ell=0\) factor is the explicitly specified integral; every
other factor is the displayed profile derivative. For example its
absolute integral is at most
\(\int_{x_-}^{x_+}\xi|\partial_\eta^sG_{J_\theta}(\xi,\eta)|\,d\xi\).
This supplies finite computable profile costs on the exact support.
The bound \(B\ge1-2h>0\), together with \(x\ge x_->0\), proves
that each fixed recursively specified coefficient is bounded on
\([x_-,x_+]\times[-1,1]\). The original moments, spacing, \(h\),
\(c_{\rm patch}\), and \(\lambda\) remain in those coefficients.

\subsection{Explicit costs in full target derivative arrays}

For \(I=(m,n)\), put \(I_z=(m+1,n)\). Use the original actual-shape
arrays \(\mathcal B_{e,I}\) from \eqref{mcq:B}, with their two
retained coupled amplitude components and old-wave cross terms.
Define the full target majorants
\begin{equation}\label{mcmean:target_arrays}
\begin{split}
 \mathsf P_I&=|\partial_z^m\partial_t^nP_{\rm old}|
          +\mathcal B_{-1,I}+\tfrac12\mathcal B_{0,I_z},\\
 \mathsf J_{\theta,I}
      &=|\partial_z^m\partial_t^n(J_\theta)_{\rm old}|
          +\tfrac12\mathcal B_{2,I},\\
 \mathsf J_{z,I}
      &=|\partial_z^m\partial_t^n(J_z)_{\rm old}|
          +\mathcal B_{1,I}+\tfrac14\mathcal B_{2,I_z}.
\end{split}
\end{equation}
The triangle inequality in \eqref{mcmean:full_targets} and the
proved increment bounds \eqref{mcq:defect_bounds} give
\(|\partial^I\mathcal P|\le\mathsf P_I\),
\(|\partial^I\mathcal J_\theta|\le\mathsf J_{\theta,I}\), and
\(|\partial^I\mathcal J_z|\le\mathsf J_{z,I}\).
Thus the old absolute derivatives appear explicitly in every cost.

For a nonnegative target array \(\mathsf X\), define the following
pointwise, fully specified finite sum:
\begin{equation}\label{mcmean:cost_operator}
\begin{split}
 \mathfrak C_{\sigma,F;k,m,n}[\mathsf X](r,z,t)
  =\sum_{i=0}^m\sum_{j=0}^n&\binom mi\binom nj\mathsf X_{(i,j)}(z,t)\\
 &\cdot q^{\sigma-k/2-(m-i)D-(n-j)}
       |F_{\sigma;k,m-i,n-j}(x,\eta)|.
\end{split}
\end{equation}
All powers are powers of the actual positive \(q\), even when their
exponents are negative. Define
\begin{equation}\label{mcmean:velocity_arrays}
\begin{split}
 \mathsf A_{\theta;k,m,n}
    &=\mathfrak C_{A,V_P;k,m,n}[\mathsf P]
        +\mathfrak C_{A-1,V_{J_z};k,m,n}[\mathsf J_z],\\
 \mathsf A_{z;k,m,n}
    &=\mathfrak C_{A-3/2,G_{J_\theta};k,m,n}[\mathsf J_\theta],\\
 \mathsf S_{k,m,n}
    &=\mathfrak C_{A-1,H_\theta;k,m,n}[\mathsf J_\theta],\\
 \mathsf A_{r;k,m,n}
    &=\mathfrak C_{A-1,H_\theta;k,m,n}[\mathsf J_\theta^{+z}]
      +\mathfrak C_{A-1-D,\mathcal Z_{A-1}H_\theta;k,m,n}
                                                 [\mathsf J_\theta],
 \qquad (\mathsf J_\theta^{+z})_{(i,j)}=\mathsf J_{\theta,(i+1,j)}.
\end{split}
\end{equation}
Equations \eqref{mcmean:all_derivatives}--\eqref{mcmean:potential_radial}
and the triangle inequality prove, for every \(k,m,n\),
\begin{equation}\label{mcmean:velocity_bounds}
\begin{gathered}
 |\partial_r^k\partial_z^m\partial_t^n\delta v|
       \le\mathsf A_{\theta;k,m,n},\qquad
 |\partial_r^k\partial_z^m\partial_t^n\delta\gamma|
       \le\mathsf A_{z;k,m,n},\\
 |\partial_r^k\partial_z^m\partial_t^n\Psi|
       \le\mathsf S_{k,m,n},\qquad
 |\partial_r^k\partial_z^m\partial_t^n\delta\beta|
       \le\mathsf A_{r;k,m,n}.
\end{gathered}
\end{equation}
The alternative upper bound \(\mathsf S_{k,m+1,n}\) for the last
quantity follows directly from \eqref{mcmean:radial_all_derivatives}.
For example the zero-order costs read
\[
\begin{split}
 |\delta v|&\le q^A\mathsf P_{0,0}|V_P|
                    +q^{A-1}\mathsf J_{z,0,0}|V_{J_z}|,\\
 |\delta\gamma|&\le q^{A-3/2}\mathsf J_{\theta,0,0}|G_{J_\theta}|,
 \qquad |\Psi|\le q^{A-1}\mathsf J_{\theta,0,0}|H_\theta|,\\
 |\delta\beta|&\le q^{A-1}\mathsf J_{\theta,1,0}|H_\theta|
       +q^{A-1-D}\mathsf J_{\theta,0,0}|\mathcal Z_{A-1}H_\theta|.
\end{split}
\]
These display separately the target derivative and moving-scale radial
cost. Replacing each profile value in \eqref{mcmean:cost_operator}
by its maximum over its fixed \(x,\eta\) support supplies a pointwise
bound depending only on \(z,t\) and the stated fixed profile data.

The same formula gives weighted radial costs without a loss hidden in
a constant. For any real \(e\), replace the final profile absolute
value in \eqref{mcmean:cost_operator} by
\[
 q^{(e+1)/4}
 \left(\int_{x_-}^{x_+}x^e
          |F_{\sigma;k,m-i,n-j}(x,\eta)|^2\,dx\right)^{1/2}.
\]
Call the resulting sum \(\mathfrak C^{[e]}\). The exact measure
identity \(r^e\,dr=q^{(e+1)/2}x^e\,dx\), followed by the triangle
inequality in this weighted \(L^2\) space, proves all four bounds
\eqref{mcmean:velocity_bounds} with the left sides replaced by their
\([e]\) norms and the right sides by the corresponding
\(\mathfrak C^{[e]}\) sums. The support is bounded away from zero,
so these integrals are finite for every fixed real \(e\).

\subsection{Actual pressure reconstruction and the surviving defects}

Set \(\boldsymbol a=(a_r,a_\theta,a_z)
=(\delta\beta,\delta v,\delta\gamma)\). Before pressure is
changed, let \(\delta g_r\) denote the exact change in the source's
angularly averaged negative radial residual. The angular average here
retains auxiliary dependence. The actual pressure-defect change is
\begin{equation}\label{mcmean:pressure_defect_step}
 \delta P_{\rm step}=\int_0^\infty\overline{\delta g_r}\,dr
                    =P_{\rm new}-\mathcal P.
\end{equation}
This is the defect change caused by the present mean operation;
it is different from \(\Delta P\) in \eqref{mcmean:full_targets}.
The source pressure bump is still
\(\rho=q^{-1/2}\widehat\rho(x)\),
\(\int\widehat\rho\,dx=1\). The actual pressure increment for this
operation is
\begin{equation}\label{mcmean:actual_pressure}
 \delta p=\mathcal I_c f,\qquad
 f=\delta g_r-\rho\,\delta P_{\rm step}.
\end{equation}
The original physical shifted primitive is
\[
\begin{split}
 \mathcal If(r,Y)&=\int_0^r
  f(s,Y+(s^{d_r}-r^{d_r})v_r,z,t)\,ds,\\
 \mathcal Jf(r,Y)&=\int_0^\infty
  f(s,Y+(s^{d_r}-r^{d_r})v_r,z,t)\,ds,\\
 \mathcal I_cf&=\mathcal If-\chi_m\mathcal Jf,\\
 d_r&=2((1+h)\rho_g-h\kappa_s),\quad
 \rho_g=\frac{\log(4-\sqrt2)}{\log(4+\sqrt2)},\quad
 \kappa_s=10^{-5},\quad v_r=(1,-(\sqrt2-1)).
\end{split}
\]
Here \(\chi_m\) is the actual inner-zero, outer-one source cutoff;
the two suppressed arguments of each integrand are exactly \(z,t\).
The pressure entering the physical force is the physical evaluation of
this source function, with its original phase evaluation and all its
derivatives. Thus it includes any nonzero auxiliary modes of the
source reconstruction.

To verify the reconstruction identity, the source physical radial
operator \(\mathcal D_r\) cancels the derivative of the term
\(-r^{d_r}v_r\) in the shift. Endpoint differentiation gives
\(\mathcal D_r\mathcal If=f\) and
\(\mathcal D_r\mathcal Jf=0\), and hence
\begin{equation}\label{mcmean:pressure_remainder}
 \mathcal D_r\delta p-\delta g_r
 =-\rho\,\delta P_{\rm step}
                  -(\partial_r\chi_m)\mathcal Jf.
\end{equation}
Haar translation invariance and \eqref{mcmean:pressure_defect_step}
give \(\overline{\mathcal Jf}=0\), but this does not set
\(\mathcal Jf\) itself to zero. In the inner collar \(\mathcal If=0\)
and \(\chi_m=0\); in the outer collar \(\mathcal If=\mathcal Jf\)
and \(\chi_m=1\). This proves compact support of the actual pressure
and all derivatives. In particular its gradient is retained in the
force calculation below.

For completeness the exact remaining source terms can be written
entirely in the original current mean fields. With
\(\Delta_0=\mathcal D_r^2+r^{-1}\mathcal D_r+\partial_z^2\),
direct expansion of the radial conservative source gives
\begin{equation}\label{mcmean:Rg}
\begin{split}
 R_g:={}&\delta g_r-\frac{2V}{r}\delta v\\
 ={}&-\partial_t\delta\beta
 -(\mathcal D_r+r^{-1})
       \bigl(2b\delta\beta+2\beta\delta\beta+(\delta\beta)^2\bigr)\\
 &-\partial_z\bigl(b\delta\gamma+G\delta\beta+
       \beta\delta\gamma+\gamma\delta\beta+
       \delta\beta\delta\gamma\bigr)\\
 &+r^{-1}\bigl(2v\delta v+(\delta v)^2\bigr)
       +\nu_{\rm NS}(\Delta_0-r^{-2})\delta\beta.
\end{split}
\end{equation}
Indeed the radial conservative quadratic expression before the change
is \(-(\mathcal D_r+r^{-1})(b+\beta)^2
-\partial_z((b+\beta)(G+\gamma))+(V+v)^2/r\), together with
the wave covariance, time term, and viscosity. Subtract it from its
value after the mean change. The unchanged wave covariance cancels;
each remaining quadratic product expands to the terms displayed.
The term \(2V\delta v/r\) has been written separately in the
definition of \(R_g\). The ordinary time derivative on the increment
is exact because it is auxiliary independent.

The five physical rows then give the exact new defects
\begin{equation}\label{mcmean:nonlinear_defects}
\begin{split}
 P_{\rm new}&=\int\overline{R_g}\,dr,\\
 (J_\theta)_{\rm new}
   &=\int r^2\overline{\gamma\delta v+v\delta\gamma+
                                    \delta\gamma\delta v}\,dr,\\
 (J_z)_{\rm new}
   &=\int r\overline{2\gamma\delta\gamma+(\delta\gamma)^2}\,dr
                        -\tfrac12\int r^2\overline{R_g}\,dr.
\end{split}
\end{equation}
For the first line add \(\mathcal P\) to
\(\int(2V\delta v/r+\overline{R_g})\,dr\) and use the third
row. For the second, expanding \((G+\gamma+\delta\gamma)
(V+v+\delta v)-(G+\gamma)(V+v)\) gives base-linear terms
\(G\delta v+V\delta\gamma\), cancelled by the fourth row, and
exactly the three remaining terms displayed. For the third, expand
\((G+\gamma+\delta\gamma)^2-(G+\gamma)^2\) and the radial-source
moment \(-\tfrac12\int r^2\delta g_r\,dr\); the base-linear
terms are the fifth row. This proves all three equalities.
Neither their right sides nor the cutoff term in
\eqref{mcmean:pressure_remainder} has been removed.

\subsection{Exact full cylindrical physical force increment}

For the actual physical fields, put
\(\mathcal R(U,p)=\partial_tU+(U\cdot\nabla)U
-\nu_{\rm NS}\Delta U+\nabla p\). Ordinary derivatives in this
formula are derivatives after the physical source evaluation; they
therefore include all phase, frame, cutoff and scale derivatives of
the actual full \(U\). Expanding the two quadratic products proves
the vector identity
\begin{equation}\label{mcmean:vector_force}
 \delta\mathcal R=\partial_t\boldsymbol a
 +(U\cdot\nabla)\boldsymbol a+(\boldsymbol a\cdot\nabla)U
 +(\boldsymbol a\cdot\nabla)\boldsymbol a
 -\nu_{\rm NS}\Delta\boldsymbol a+\nabla\delta p.
\end{equation}
The pressure here is precisely \eqref{mcmean:actual_pressure}.

To derive every cylindrical component, the basis derivatives are
\(\partial_\theta e_r=e_\theta\),
\(\partial_\theta e_\theta=-e_r\), and
\(\partial_\theta e_z=0\). Thus for two arbitrary vector fields
\(X,Y\) the component vector of \((X\cdot\nabla)Y\) is
\[
 X_r\partial_rY+\frac{X_\theta}{r}\partial_\theta Y
 +X_z\partial_zY+\frac{X_\theta}{r}(-Y_\theta,Y_r,0).
\]
Apply this equality to \((U,\boldsymbol a)\),
\((\boldsymbol a,U)\), and \((\boldsymbol a,\boldsymbol a)\).
The scalar components of \(\boldsymbol a\) have zero angular
derivative, whereas those of the actual \(U\) generally do not.
Writing \(L_0=\partial_r^2+r^{-1}\partial_r+\partial_z^2\), the
complete result is
\begin{equation}\label{mcmean:force_r}
\begin{split}
 \delta\mathcal R_r={}&\partial_ta_r+U_r\partial_ra_r+U_z\partial_za_r
       +a_r\partial_rU_r+a_z\partial_zU_r
       +\frac{a_\theta}{r}\partial_\theta U_r\\
 &+a_r\partial_ra_r+a_z\partial_za_r
       -\frac{2U_\theta a_\theta+a_\theta^2}{r}
       -\nu_{\rm NS}(L_0-r^{-2})a_r+\partial_r\delta p,
\end{split}
\end{equation}
\begin{equation}\label{mcmean:force_theta}
\begin{split}
 \delta\mathcal R_\theta={}&\partial_ta_\theta
       +U_r\partial_ra_\theta+U_z\partial_za_\theta
       +a_r\partial_rU_\theta+a_z\partial_zU_\theta
       +\frac{a_\theta}{r}\partial_\theta U_\theta\\
 &+a_r\partial_ra_\theta+a_z\partial_za_\theta
       +\frac{U_ra_\theta+a_rU_\theta+a_ra_\theta}{r}
       -\nu_{\rm NS}(L_0-r^{-2})a_\theta
       +\frac1r\partial_\theta\delta p,
\end{split}
\end{equation}
\begin{equation}\label{mcmean:force_z}
\begin{split}
 \delta\mathcal R_z={}&\partial_ta_z+U_r\partial_ra_z+U_z\partial_za_z
       +a_r\partial_rU_z+a_z\partial_zU_z
       +\frac{a_\theta}{r}\partial_\theta U_z\\
 &+a_r\partial_ra_z+a_z\partial_za_z
       -\nu_{\rm NS}L_0a_z+\partial_z\delta p.
\end{split}
\end{equation}
The self terms in the radial equation include
\(-a_\theta^2/r\); those in the angular equation include
\(a_ra_\theta/r\). The old-wave cross terms occur in the actual
\(U\) and in its angular derivatives in all three equations.
The mean pressure is angularly invariant, so its angular derivative
is zero for this source operation; its radial and axial derivatives
are the actual reconstructed ones and are generally nonzero.

For the viscosity formula, applying the scalar cylindrical Laplacian
to the moving basis yields, for a general vector field \(X\),
\[
 \Delta X=
 \left(\Delta_sX_r-r^{-2}X_r-2r^{-2}\partial_\theta X_\theta,
       \Delta_sX_\theta-r^{-2}X_\theta+2r^{-2}\partial_\theta X_r,
       \Delta_sX_z\right),
 \quad\Delta_s=L_0+r^{-2}\partial_\theta^2.
\]
This follows by differentiating the two basis identities twice and
applying the product rule. Substituting the angularly independent
components of \(\boldsymbol a\) gives exactly
\(((L_0-r^{-2})a_r,(L_0-r^{-2})a_\theta,L_0a_z)\), proving all
three displayed viscous terms.
The radial expression before the pressure term, when angularly
averaged in the extended source domain and negated, is precisely
\(\delta g_r\): this follows from the definition of the negative
radial source residual and linearity of the average. It is also an
explicit formula for the input to \eqref{mcmean:actual_pressure} in
terms of the same full velocity and mean increments.

\subsection{Finite all-order force costs with every inverse-radius term}

Here are explicit derivative bounds for all three components, including
angular derivatives of the actual current velocity. Let
\(\alpha=(k,\ell,m,n)\in\mathbb N^4\) index the scalar component
derivative \(\partial^\alpha=\partial_r^k\partial_\theta^\ell
\partial_z^m\partial_t^n\), and let
\(e_r,e_\theta,e_z,e_t\) be the four coordinate multi-indices.
Extend the arrays in \eqref{mcmean:velocity_arrays} by
\[
 \mathsf A_{i,\alpha}=
 \begin{cases}\mathsf A_{i;k,m,n}&\ell=0,\\0&\ell>0,
 \end{cases}
 \qquad
 \mathsf U_{i,\alpha}=|\partial^\alpha U_i|.
\]
The pressure array \(\mathsf\Pi\) will be the following explicit
majorant for the derivatives of the actual specified function
\eqref{mcmean:actual_pressure}, after physical evaluation.

The original physical radial phase satisfies
\(Y_{\rm phys}(s,z,t)-Y_{\rm phys}(r,z,t)
=(s^{d_r}-r^{d_r})v_r\). Substituting this exact equality in the
two shifted integrals proves the evaluated formula
\begin{equation}\label{mcmean:evaluated_pressure}
\begin{split}
 \delta p(r,z,t)
 &=\int_0^r f_{\rm phys}(s,z,t)\,ds
        -\chi_m(r,z,t)\int_0^\infty f_{\rm phys}(s,z,t)\,ds,\\
 f_{\rm phys}&=(\delta g_r)_{\rm phys}
                       -\rho\,\delta P_{\rm step}.
\end{split}
\end{equation}
Every original axial and temporal phase dependence is retained in
\(f_{\rm phys}\). The second integral is independent of \(r\)
after this evaluation. This equality does not assert that it is zero.

Fix any compact interior physical domain for the finite operation,
and choose one physical interval \([R_-,R_+]\Subset(0,\infty)\)
containing the supports of \(\delta g_r\) and \(\rho\) for that
domain and every auxiliary phase. Put \(L_R=R_+-R_-\), with these
actual endpoints kept fixed during differentiation. The smooth zero
extensions imply that all their derivatives have support in the
same interval. Define \(\mathsf G_{k,m,n}(z,t)\) as the supremum
over \(r\in[R_-,R_+]\) and the original auxiliary phases of the
angular average of the following nonnegative expression: the radial
cost \(\mathsf F_{r;(k,0,m,n)}\) in \eqref{mcmean:force_costs}
with its last pressure term deleted. This definition involves only
\(\mathsf A\), the actual current velocity derivatives,
\(\nu_{\rm NS}\), and the inverse-radius arrays; it has no
dependence on \(\mathsf\Pi\). When auxiliary phases are retained,
the velocity derivative arrays mean the corresponding full lifted
physical derivatives before evaluation. The chain rule for source
evaluation identifies these with derivatives of the evaluated fields
at each phase. The force identities and the derivative proof below
therefore imply
\[
 |\partial_r^k\partial_z^m\partial_t^n(\delta g_r)_{\rm phys}|
       \le\mathsf G_{k,m,n},\qquad
 |\partial_z^m\partial_t^n\delta P_{\rm step}|
       \le\mathsf D_{m,n}:=L_R\mathsf G_{0,m,n}.
\]
For the first inequality use the negative angularly averaged radial
force before pressure, as proved above, and the triangle inequality
under that average. For the second, differentiate
\eqref{mcmean:pressure_defect_step} under its compact radial
integral, commute axial and temporal differentiation with Haar
average, and integrate the same bound over the interval of width
\(L_R\). The Haar integral of each original constant torus
direction derivative is zero by periodicity, which justifies that
commutation for the full physical time derivative as well.

The derivative of the pressure bump is already an explicit original-
scale coefficient in \eqref{mcq:bump_derivatives}. Define its exact
support supremum
\[
 \mathsf H_{k,m,n}(z,t)
 =q^{-(1+k)/2-mD-n}
    \sup_{r\in[R_-,R_+]}
       |f_{0;k,m,n}(r/\sqrt q,\eta)|,
\]
where \(f_{0;k,m,n}\) is the recursion starting with the actual
\(\widehat\rho\), rather than another normalized bump.
It follows that \(|\partial_r^k\partial_z^m\partial_t^n\rho|
\le\mathsf H_{k,m,n}\). Set
\begin{equation}\label{mcmean:pressure_source_cost}
 \mathsf E_{k,m,n}
 =\mathsf G_{k,m,n}
 +\sum_{i=0}^m\sum_{j=0}^n\binom mi\binom nj
        \mathsf H_{k,m-i,n-j}\mathsf D_{i,j}.
\end{equation}
Since \(\delta P_{\rm step}\) has no radial dependence, all
\(k\) radial derivatives in its product with \(\rho\) act on
\(\rho\). The repeated axial and temporal product rule proves
\(|\partial_r^k\partial_z^m\partial_t^n f_{\rm phys}|
\le\mathsf E_{k,m,n}\) on the full interval.

Write \(\mathsf X_{k,i,j}(r,z,t)
=|\partial_r^k\partial_z^i\partial_t^j\chi_m(r,z,t)|\), retaining
the actual cutoff and its moving-scale derivatives. Define
\begin{equation}\label{mcmean:explicit_pressure_cost}
\begin{split}
 \mathsf\Pi_{(k,0,m,n)}={}&
 \begin{cases}L_R\mathsf E_{0,m,n}&k=0,\\
               \mathsf E_{k-1,m,n}&k\ge1
 \end{cases}\\
 &+L_R\sum_{i=0}^m\sum_{j=0}^n\binom mi\binom nj
           \mathsf X_{k,i,j}\mathsf E_{0,m-i,n-j},
 \qquad \mathsf\Pi_{(k,\ell,m,n)}=0\quad(\ell>0).
\end{split}
\end{equation}
This proves the explicit bound
\(|\partial^\alpha\delta p|\le\mathsf\Pi_\alpha\).
In fact differentiating \eqref{mcmean:evaluated_pressure} yields
the exact first-integral derivative
\[
 \partial_r^k\partial_z^m\partial_t^n
       \int_0^r f_{\rm phys}(s,z,t)\,ds
 =\begin{cases}
   \displaystyle\int_0^r\partial_z^m\partial_t^n
                         f_{\rm phys}(s,z,t)\,ds&k=0,\\
   \partial_r^{k-1}\partial_z^m\partial_t^n f_{\rm phys}(r,z,t)&k\ge1.
  \end{cases}
\]
The absolute values are bounded by the first term in
\eqref{mcmean:explicit_pressure_cost}. In the cutoff product all
radial derivatives act on \(\chi_m\), whereas the axial and temporal
derivatives distribute with the two displayed binomial factors.
Every derivative of its radial integral is bounded by
\(L_R\mathsf E_{0,m-i,n-j}\). This proves its second term, with
every cutoff derivative and the actual support width retained.
Angular invariance proves the stated zero entries. This construction
first computes the pressure-free \(\mathsf G\), then
\(\mathsf D,\mathsf H,\mathsf E,\mathsf\Pi\), and finally the
full force array below, so the pressure budget is a finite calculation
without circularly assuming a pressure bound.

For nonnegative arrays define their exact Leibniz convolution and
derivative shift by
\begin{equation}\label{mcmean:convolution}
 (X\star Y)_\alpha=\sum_{\beta\le\alpha}
               \binom\alpha\beta X_\beta Y_{\alpha-\beta},
 \qquad (S_\delta X)_\alpha=X_{\alpha+\delta}.
\end{equation}
Products with three factors use either parenthesization; multiplying
the binomial coefficients gives
\(\alpha!/(\beta!\gamma!(\alpha-\beta-\gamma)!)\), proving the
two parenthesizations agree. For the exact inverse-radius arrays put
\begin{equation}\label{mcmean:radius_array}
 (\mathsf R_j)_\alpha=
 \begin{cases}
 (j)_k r^{-j-k}&\ell=m=n=0,\\0&\text{otherwise},
 \end{cases}
 \quad (j)_0=1,\quad(j)_k=j(j+1)\cdots(j+k-1),\quad j=1,2.
\end{equation}
Induction gives \(\partial_r^k r^{-j}=(-1)^k(j)_kr^{-j-k}\),
which proves that \(\mathsf R_j\) is its exact absolute derivative
array. No radial derivative of an inverse-radius coefficient is
omitted in the following costs.

For \(i=r,\theta,z\) define the common transport cost
\begin{equation}\label{mcmean:transport_cost}
\begin{split}
 \mathsf T_i={}&S_{e_t}\mathsf A_i
 +\mathsf U_r\star S_{e_r}\mathsf A_i
 +\mathsf U_z\star S_{e_z}\mathsf A_i
 +\mathsf A_r\star S_{e_r}\mathsf U_i
 +\mathsf A_z\star S_{e_z}\mathsf U_i\\
 &+\mathsf R_1\star\mathsf A_\theta\star S_{e_\theta}\mathsf U_i
 +\mathsf A_r\star S_{e_r}\mathsf A_i
 +\mathsf A_z\star S_{e_z}\mathsf A_i,
\end{split}
\end{equation}
and the exact finite upper-cost arrays
\begin{equation}\label{mcmean:force_costs}
\begin{split}
 \mathsf F_r={}&\mathsf T_r
    +2\mathsf R_1\star\mathsf U_\theta\star\mathsf A_\theta
    +\mathsf R_1\star\mathsf A_\theta\star\mathsf A_\theta\\
 &+\nu_{\rm NS}\bigl(S_{2e_r}\mathsf A_r
       +\mathsf R_1\star S_{e_r}\mathsf A_r
       +S_{2e_z}\mathsf A_r+\mathsf R_2\star\mathsf A_r\bigr)
       +S_{e_r}\mathsf\Pi,\\
 \mathsf F_\theta={}&\mathsf T_\theta
    +\mathsf R_1\star\bigl(\mathsf U_r\star\mathsf A_\theta
                    +\mathsf A_r\star\mathsf U_\theta
                    +\mathsf A_r\star\mathsf A_\theta\bigr)\\
 &+\nu_{\rm NS}\bigl(S_{2e_r}\mathsf A_\theta
       +\mathsf R_1\star S_{e_r}\mathsf A_\theta
       +S_{2e_z}\mathsf A_\theta+\mathsf R_2\star\mathsf A_\theta\bigr)
       +\mathsf R_1\star S_{e_\theta}\mathsf\Pi,\\
 \mathsf F_z={}&\mathsf T_z
    +\nu_{\rm NS}\bigl(S_{2e_r}\mathsf A_z
       +\mathsf R_1\star S_{e_r}\mathsf A_z
       +S_{2e_z}\mathsf A_z\bigr)+S_{e_z}\mathsf\Pi.
\end{split}
\end{equation}
Every addition and scalar multiplication here is entrywise. For every
finite \(\alpha\), these arrays prove the full physical component
inequalities
\begin{equation}\label{mcmean:all_force_bounds}
 |\partial^\alpha\delta\mathcal R_i|\le\mathsf F_{i,\alpha},
 \qquad i=r,\theta,z.
\end{equation}
Indeed the repeated Leibniz rule gives
\(|\partial^\alpha(fg)|\le(X\star Y)_\alpha\) whenever the
respective derivative absolute values are bounded by \(X,Y\).
The three-factor version is its iterated rule, proved by the explicit
multinomial coefficient above. Apply these rules term by term to
\eqref{mcmean:force_r}--\eqref{mcmean:force_z}, using
\eqref{mcmean:velocity_bounds} for the increment, the actual
derivatives for \(U,\delta p\), and \eqref{mcmean:radius_array} for
each inverse radius. The constant viscosity passes through every
derivative. The three resulting sums are precisely
\eqref{mcmean:transport_cost}--\eqref{mcmean:force_costs}, proving
the claim at every order.

These are component derivatives in the original orthonormal basis.
For an explicit bound on the actual Cartesian vector after angular
differentiation, the moving basis gives
\[
 \left|\partial_r^k\partial_\theta^\ell\partial_z^m\partial_t^n
                 \delta\mathcal R\right|
 \le\sum_{j=0}^\ell\binom\ell j
       \bigl(\mathsf F_{r;(k,j,m,n)}
             +\mathsf F_{\theta;(k,j,m,n)}\bigr)
       +\mathsf F_{z;(k,\ell,m,n)}.
\]
To prove it, expand each angular derivative by the binomial rule
between a scalar component and its basis vector. Every derivative of
\(e_r,e_\theta\) has norm one, whereas the derivatives of \(e_z\)
of positive order vanish. The triangle inequality gives exactly the
displayed bound. Thus the component formulas also provide a bound
for the physical vector, with its angular basis differentiation.

There is also an exact recursion for every Cartesian derivative.
Let \(O_\theta\) be the orthogonal matrix with columns
\((e_r,e_\theta,e_z)\), and let
\(J(f_r,f_\theta,f_z)=(-f_\theta,f_r,0)\). Use
\(x_{\rm C}=r\cos\theta,y_{\rm C}=r\sin\theta\) for the physical
Cartesian coordinates, distinct from the original scaled radius
\(x=r/\sqrt q\). Define the operators on component vectors
\begin{equation}\label{mcmean:cartesian_recursion}
\begin{split}
 \mathcal Xf&=\cos\theta\,\partial_rf
        -\frac{\sin\theta}{r}(\partial_\theta f+Jf),\\
 \mathcal Yf&=\sin\theta\,\partial_rf
        +\frac{\cos\theta}{r}(\partial_\theta f+Jf),
 \qquad \mathcal Zf=\partial_zf.
\end{split}
\end{equation}
The scalar coordinate rules
\(\partial_{x_{\rm C}}=\cos\theta\partial_r
-r^{-1}\sin\theta\partial_\theta\) and
\(\partial_{y_{\rm C}}=\sin\theta\partial_r
+r^{-1}\cos\theta\partial_\theta\), together with
\(\partial_\theta O_\theta=O_\theta J\), give exactly
\(\partial_{x_{\rm C}}(O_\theta f)=O_\theta\mathcal Xf\),
\(\partial_{y_{\rm C}}(O_\theta f)=O_\theta\mathcal Yf\), and
\(\partial_z(O_\theta f)=O_\theta\mathcal Zf\).
Applying these identities repeatedly proves
\[
 \partial_{x_{\rm C}}^i\partial_{y_{\rm C}}^j
 \partial_z^k\partial_t^n(O_\theta f)
   =O_\theta\mathcal X^i\mathcal Y^j\mathcal Z^k\partial_t^nf.
\]
The order written is a fixed composition order and is exact because
it was obtained from commuting smooth Cartesian derivatives.
For a fully explicit majorant recursion, if \(\mathsf C_r,
\mathsf C_\theta,\mathsf C_z\) majorize the component derivative
arrays of \(f\), set \((\mathsf J C)_r=\mathsf C_\theta\),
\((\mathsf J C)_\theta=\mathsf C_r\),
\((\mathsf J C)_z=0\). Let \(\mathsf c_\alpha
=|\partial^\alpha\cos\theta|\),
\(\mathsf s_\alpha=|\partial^\alpha\sin\theta|\), whose entries
vanish unless \(k=m=n=0\) and whose remaining entries are the
actual derivatives of these elementary functions. Then one
\(\mathcal X\) application is bounded componentwise by
\[
 \mathsf c\star S_{e_r}\mathsf C_i
 +\mathsf s\star\mathsf R_1\star
             (S_{e_\theta}\mathsf C_i+(\mathsf J C)_i),
\]
and one \(\mathcal Y\) application by the same expression with
\(\mathsf c,\mathsf s\) interchanged. A \(\mathcal Z\) or time
application is the array shift \(S_{e_z}\) or \(S_{e_t}\).
These rules follow from the exact product rule
\eqref{mcmean:convolution}. Starting with \(\mathsf F\), iterating
them the indicated finite number of times, and taking the sum of
the three resulting zero-order component bounds gives a finite bound
for each stated Cartesian vector derivative. Thus all Cartesian
inverse-radius and angular-basis costs are also explicitly supplied.

All costs above are finite for the actual finite state on each compact
interior physical domain. Their dependence on the original scale,
target derivatives, inverse-radius factors, viscosity, full current
velocity, and reconstructed pressure is explicit. The exact nonlinear
defects \eqref{mcmean:nonlinear_defects} remain after this operation.
No uniform estimate over an increasing family of source stages,
no singular-time limit, and no infinite correction gain is asserted
by these finite derivative identities and costs.

\end{document}